

\documentclass[a4paper,11pt]{article}

\usepackage[utf8]{inputenc}
\usepackage{amsmath,amssymb,amsthm,mathtools}
\usepackage{graphicx}
\usepackage{booktabs}
\usepackage{hyperref}
\usepackage{cite}
\usepackage{geometry}
\usepackage{bm}
\usepackage{lmodern}
\usepackage{microtype}

\geometry{margin=1in}

\newtheorem{theorem}{Theorem}[section]
\newtheorem{proposition}[theorem]{Proposition}
\newtheorem{lemma}[theorem]{Lemma}
\newtheorem{corollary}[theorem]{Corollary}
\theoremstyle{definition}
\newtheorem{definition}[theorem]{Definition}
\newtheorem{remark}[theorem]{Remark}
\newcommand{\Tr}{\mathop{\rm Tr}}
\newcommand{\MP}{M_{\rm P}}
\newcommand{\Lobs}{\Lambda_{\rm obs}}
\newcommand{\Icl}{\mathcal{I}_{\rm cl}}
\newcommand{\Ilat}{\mathcal{I}_{\rm lat}}
\newcommand{\Jphys}{J_{\rm phys}}
\newcommand{\bR}{\mathbb{R}}
\newcommand{\Bobs}{\mathcal{B}}

\title{A lattice computation of the nonperturbative cosmological-constant residual\\
from $SU(2)$ gravitational-instanton branch bias}

\author{M.~M.~Hanash\thanks{ORCID: 0009-0002-8426-825X. Independent researcher.}}
\date{\today}

\begin{document}
\maketitle

\begin{abstract}
We compute the cosmological constant on the lattice within the
Sakharov--Volovik induced-gravity tradition, taking as framework
input a coframe / teleparallel completion of the gravitational
effective action and an analytic constitutive subtraction that
removes polynomial ultraviolet hierarchies order by order in the
one-loop branch-bias expansion. The surviving infrared residual
is reopened nonperturbatively by a self-dual $SU(2)$ gravitational
instanton of Harrington--Shepard type on $S^1 \times \bR^3$,
whose action supplies a dimensional-transmutation exponential and
whose one-loop branch determinant supplies an $\mathcal{O}(1)$
prefactor. We compute on three independent lattices spanning a
factor of $11$ in volume and a factor of $1.5$ in caloron scale
the three lattice ingredients entering that prefactor: the
physical Jacobian of collective coordinates, the phase of the
instanton contribution, and the normalized branch-response
bracket together with its finite-$\xi$ Simpson correction and
lattice-size controls. The branch-response determinant is
evaluated with the adaptive-prime prescription validated
internally by lattice-size stability and by an independent
projected-complement cross-check. We establish a factorization
proposition for the master formula, a phase theorem fixing the
parity-even angle to $+1$ on every background considered, and a
stability proposition for the Simpson nonlinear correction across
the three production lattices. Conditional on the hidden-sector action $S_{\rm inst} \simeq 276$
supplied by a companion analysis, the assembled
lattice-determined prefactor is $0.81 \pm 0.03$, so that
\[
\Lobs/\MP^4 \;=\; (0.81 \pm 0.03) \times 10^{-120}
\qquad
\text{(conditional on $S_{\rm inst} \simeq 276$)}.
\]
\end{abstract}

\section{Introduction}
\label{sec:introduction}

The cosmological constant problem is the discrepancy between the
small positive vacuum energy density inferred from cosmological
observations and the much larger values expected from effective
field theory applied to the known matter sectors of the Standard
Model and its
extensions~\cite{Weinberg89,Carroll01,Padmanabhan03,Martin12}.
Reproducing the observed magnitude
$\Lobs \sim 10^{-120}\,\MP^4$ within a controlled theoretical
framework, without relying on anthropic selection or on
cancellations between unrelated sectors, has remained an open
problem in high-energy theory. Several distinct approaches have
been proposed, including supersymmetric cancellation, adjustment
mechanisms based on scalar or four-form sectors, sequestering
constructions, and vacuum-selection arguments; a comparison of
the present paper to these approaches is given in
Section~\ref{sec:discussion}.

We work within the Sakharov--Volovik induced-gravity
tradition~\cite{Sakharov67,Visser02,Volovik03,Volovik2001}, in
which the gravitational effective action is generated by
integrating out fluctuations of a microscopic proto-field sector
on an emergent metric $g_{\rm eff}$. The Newton constant and the
cosmological constant are both induced by this procedure, so that
their magnitudes are fixed by the microscopic content of the
proto-field sector and by whatever cancellation or subtraction
structure the framework supplies. In the specific framework of
the companion structural paper~\cite{companion-structural}%
\footnote{Available at DOI \href{https://doi.org/10.5281/zenodo.19599312}{10.5281/zenodo.19599312}; under journal review.}%
, the
induced gravitational action is completed by a coframe /
teleparallel sector that selects a preferred FLRW foliation, and
polynomial ultraviolet hierarchies in the vacuum-energy response
between two constitutive branches are removed by an analytic
subtraction within that sector. The result of the subtraction is
a finite infrared residual, which the present paper computes. The companion paper is deposited at DOI~10.5281/zenodo.19599312 and
is under journal review.

The residual is reopened nonperturbatively by a self-dual $SU(2)$
gravitational instanton of Harrington--Shepard
type~\cite{HS1978} on $S^1 \times \bR^3$, constructed from
the underlying BPST instanton~\cite{BPST1975} with the 't~Hooft
sign conventions~\cite{tHooft1976}. The
instanton action supplies a dimensional-transmutation
exponential, while the one-loop branch determinant computed on
the instanton background supplies the normalized branch-response
factor entering the lattice prefactor. The resulting master
formula for the cosmological constant,
\begin{equation*}
\Lobs \;\sim\;
\MP^4\,e^{-8\pi^2 a_C \Ilat}\,
\Jphys\,\bigl[\zeta(0) - \zeta(X_*)\bigr]\,\cos\Theta_\gamma,
\end{equation*}
factorizes into four UV-finite quantities after the analytic
subtraction. The structural derivation of this factorization and
of the accompanying no-go results ruling out local scalar-only
realizations of the same residual is given
in~\cite{companion-structural}; the present paper takes that
derivation as input.

What this paper establishes is a lattice computation of the three
factors of the master formula that are intrinsic to the
branch-bias determinant: the physical Jacobian $\Jphys$, the
phase $\cos\Theta_\gamma$, and the normalized branch-response
bracket $\Bobs \equiv 1 - \zeta(X_*)/\zeta(0)$. The
dimensional-transmutation exponential $e^{-S_{\rm inst}}$ is not
recomputed here; it is taken as conditional input
from~\cite{companion-structural}. The physical Jacobian of the
collective-coordinate measure is computed in
Section~\ref{sec:jacobian}; the phase is determined by a theorem
within the parity-even class in Section~\ref{sec:phase}; the
branch-response bracket is measured on three independent lattices
via a transported-prime Simpson campaign in
Section~\ref{sec:finite-xi}, with lattice-size independence
established in Section~\ref{sec:continuum}. Three named results
form the logical spine of the computation: a factorization
proposition establishing UV-finiteness of each master-formula
factor, a phase theorem fixing $\cos\Theta_\gamma = +1$ in the
parity-even class on every background geometry considered, and a
stability proposition for the Simpson nonlinear correction across
three lattices spanning a factor of $11$ in volume. Combined with
the conditional hidden-sector input $S_{\rm inst} \simeq 276$
from~\cite{companion-structural}, the lattice computation yields
an assembled prefactor $\mathcal{K} = 0.81 \pm 0.03$ multiplying
the dimensional-transmutation exponential $10^{-120}$; the
canonical assembled statement is given in
Section~\ref{sec:assembled}.

\paragraph{Organization.}
Section~\ref{sec:framework} fixes the master formula and the
framework inputs inherited from~\cite{companion-structural}.
Section~\ref{sec:lattice-setup} defines the lattice
discretization and admissibility window.
Section~\ref{sec:jacobian} computes the physical Jacobian.
Section~\ref{sec:phase} establishes the phase theorem.
Section~\ref{sec:spectral} describes the spectral structure of
the lattice vector operator and the adaptive-prime prescription
used in the determinant computation. Section~\ref{sec:history}
reviews the prescriptions considered and reports an independent
cross-prescription measurement. Section~\ref{sec:finite-xi}
presents the transported-prime Simpson campaign and the
three-row production result. Section~\ref{sec:continuum}
establishes lattice-size independence across the three production
rows. Section~\ref{sec:assembled} assembles the final result
with the full error budget. Section~\ref{sec:discussion}
discusses the conditional scope, the coupling assumption, and
the relation to existing approaches; Section~\ref{sec:conclusion}
concludes. Six appendices collect the operator and solver
conventions, the transported-prime algorithm, the
infrastructure-test suite, the $\Icl$-degraded reference rows,
the bracket error budget, and the sensitivity of the assembled
result to the hidden-sector action.


\section{Master formula and inherited framework inputs}
\label{sec:framework}

\subsection{Framework inputs}
\label{subsec:framework-inputs}

The framework used below, derived in the companion structural
paper~\cite{companion-structural} and treated here as input, has three
ingredients. First, a proto-field sector whose one-loop integration
yields an infrared gravitational effective action
$\Gamma_{\rm grav}[g_{\rm eff}]$ in the
Sakharov~\cite{Sakharov67,Visser02}--Volovik~\cite{Volovik03,Volovik2001}
induced-gravity tradition. Second, a coframe / teleparallel completion
that provides the geometric branch structure on which the constitutive
scalar $X$ is defined. Third, an analytic constitutive subtraction
that removes all polynomial ultraviolet terms in the one-loop
branch-bias expansion, leaving a finite infrared residual. None of
these ingredients is re-derived in the present paper. The purpose of
the sections that follow is to compute, on the lattice, the residual
left after this subtraction.

\subsection{Branch-bias mechanism and master formula}
\label{subsec:framework-master}

The constitutive sector admits two branches, distinguished by the
value $X = 0$ or $X = X_*$ of a constitutive scalar. After the
analytic subtraction of~\cite{companion-structural}, the difference of
the one-loop determinants between the two branches is UV-finite, and
the master formula for the observed cosmological constant is
\begin{equation}
\Lobs \;\sim\;
\MP^4\,
e^{-8\pi^2 a_C \Ilat}\,
\Jphys\,
\bigl[\zeta(0) - \zeta(X_*)\bigr]\,
\cos\Theta_\gamma .
\label{eq:master}
\end{equation}
The factors are: the dimensional-transmutation exponential
$e^{-S_{\rm inst}}$ with $S_{\rm inst} = 8\pi^2 a_C \Ilat$, where
$a_C$ is the constitutive-sector bare coupling and $\Ilat$ the
normalized lattice instanton action, both as defined
in~\cite{companion-structural}; the physical Jacobian of collective
coordinates $\Jphys$; the phase of the instanton contribution
$\cos\Theta_\gamma$; and the one-loop branch-response bracket
$\Bobs \equiv 1 - \zeta(X_*)/\zeta(0)$. Equivalently,
\[
\zeta(0) - \zeta(X_*) \;=\; \zeta(0)\,\Bobs ,
\]
so that the bracket $\Bobs$ is the normalized branch-response factor
quoted throughout the paper.

\begin{proposition}[Factorization and finiteness]
\label{prop:factor-partition}
After the analytic constitutive subtraction
of~\cite{companion-structural}, the master formula~\eqref{eq:master}
factorizes into four UV-finite quantities. In the present paper,
$\Jphys$, $\cos\Theta_\gamma$, and
$\Bobs \equiv 1 - \zeta(X_*)/\zeta(0)$ are computed directly on the
lattice, whereas $e^{-S_{\rm inst}}$ is treated as an inherited input
from the companion structural analysis.
\end{proposition}

\begin{proof}
The finiteness of the factorization is inherited from the analytic
subtraction proved in~\cite{companion-structural}. The Jacobian
$\Jphys$ is the square root of a $5\times 5$ Gram determinant of
collective-coordinate tangents and is finite by construction. The
phase $\cos\Theta_\gamma$ is fixed by the parity theorem of
Section~\ref{sec:phase}. For the branch-response factor, the same
subtraction renders $\log\zeta(X) - \log\zeta(0)$ polynomial-free
order by order in the ultraviolet expansion, leaving the finite
residual computed in Section~\ref{sec:finite-xi}. The exponential
factor $e^{-S_{\rm inst}}$ is not recomputed here; it enters as a
conditional framework input.
\end{proof}

\subsection{Conditional scope of this paper}
\label{subsec:framework-scope}

Two conditionals are fixed in the computation that follows, and the
assembled result of Section~\ref{sec:assembled} is conditional on
both. The hidden-sector action entering the dimensional-transmutation
exponential is taken as $S_{\rm inst} \simeq 276$ from a separate
analysis of the hidden sector in~\cite{companion-structural}; all
numerical content of the present paper is independent of the
hidden-sector algebra, $S_{\rm inst}$ enters the master formula only
multiplicatively, and the sensitivity of the assembled result to its
value is quantified in Appendix~\ref{app:S-sensitivity}. The
branch-response determinant is computed using the adaptive-prime
prescription of Section~\ref{sec:spectral}. In the present paper that
prescription is supported by the lattice-size stability tests of
Section~\ref{sec:continuum} and by the independent
projected-complement cross-check
of~Section~\ref{subsec:cross-check-exact}. A first-principles dressed
analytic basis for the adaptive prime is deferred to separate
technical work. Nothing in the numerical pipeline below depends on
the hidden-sector field content beyond the single number
$S_{\rm inst}$.



\section{Lattice discretization of the Harrington--Shepard caloron}
\label{sec:lattice-setup}

\subsection{Seed construction and path-ordered links}
\label{subsec:seed-construction}

The Harrington--Shepard (HS) caloron~\cite{HS1978} is the
trivial-holonomy unit-charge periodic BPST instanton on
$S^1 \times \bR^3$~\cite{BPST1975}. In the notation used here its
gauge potential is written
\begin{equation}
A_\mu^a(x)
\;=\;
\eta^{(\pm) a}_{\mu\nu}\,\partial_\nu \log\Pi(x),
\qquad
\Pi(x)
\;=\;
1 \;+\;
\frac{\pi\,\rho^2}{\beta\,r}\,
\frac{\sinh(2\pi r/\beta)}
     {\cosh(2\pi r/\beta) \;-\; \cos(2\pi t/\beta)},
\label{eq:HS-potential}
\end{equation}
where $\rho$ is the caloron scale parameter, $r = |\vec{x}|$, and
$\eta^{(\pm) a}_{\mu\nu}$ denotes the 't~Hooft symbol of the chosen
orientation~\cite{tHooft1976}. Lattice links $U_{x,\mu} \in SU(2)$
are constructed by discretized path-ordered integration
of~\eqref{eq:HS-potential} along each lattice edge into
$n_{\rm sub} = 16$ sub-intervals, composed as an ordered product of
sub-link exponentials, and unitarity-projected onto $SU(2)$ by polar
decomposition; the absolute deviation
$\|U^\dagger U - \mathbb{1}\|_{\max}$ stays below $10^{-15}$ on every
production background.

\subsection{Admissibility window and clover diagnostics}
\label{subsec:admissibility}

The quality of the discretized seed is monitored by the
clover-averaged lattice instanton action
\begin{equation}
\Icl
\;\equiv\;
\frac{1}{8\pi^2}\sum_x
\Tr\!\bigl[F^{\rm clover}_{\mu\nu}(x)\,F^{\rm clover\,\mu\nu}(x)\bigr],
\label{eq:Icl-def}
\end{equation}
where $F^{\rm clover}_{\mu\nu}$ is the symmetric average of the four
plaquettes meeting at $x$ in the $(\mu,\nu)$-plane. In the present
paper, $\Icl$ is used only as a discretization-quality diagnostic for
the seeded caloron background. The continuum target for a $Q = 1$
self-dual caloron is $\Icl = 1$; deviations measure the combined
effect of finite-spacing and finite-volume artefacts on the seeded
profile.

The target admissibility window for the primary production rows is
\begin{equation}
\Icl \;\in\; [\,0.98,\,1.00\,].
\label{eq:icl-window}
\end{equation}
In addition, one row at $\Icl = 1.032$ is retained in the main
analysis as a stress test of the window itself: its inclusion in
the weighted bracket of Section~\ref{sec:finite-xi} is designed to
test whether the bracket value is sensitive to mild excursions
from~\eqref{eq:icl-window}. The empirical answer, reported in
Section~\ref{sec:finite-xi}, is that it is not. Rows with larger
departures from unity are used only as a systematics reference and
are deferred to Appendix~\ref{app:icl-degraded}. All production
backgrounds are seeded from the unit-charge HS caloron; no separate
lattice measurement of $Q$ is used in the present paper.

Table~\ref{tab:production-rows} lists the three production rows used
in the remainder of the paper.

\begin{table}[h]
\centering
\begin{tabular}{lccc}
\toprule
Row & Dimensions & $\rho$ & $\Icl$ \\
\midrule
A & $18^3 \times 12$ & $3.0$ & $0.983$ \\
B & $24^3 \times 16$ & $3.5$ & $0.998$ \\
C & $32^3 \times 24$ & $4.5$ & $1.032$ \\
\bottomrule
\end{tabular}
\caption{Production rows. Rows A and B lie inside the target
admissibility window~\eqref{eq:icl-window}; Row C is retained as an
explicit stress test slightly outside it, extending the range to a
factor of $11$ in lattice volume and a factor of $1.5$ in caloron
scale.}
\label{tab:production-rows}
\end{table}

\subsection{Reference conventions}
\label{subsec:reference-conventions}

Throughout the production pipeline the infrared regulator is
$m^2 = 0.01$ and the disformal branch parameter is $\epsilon = 0.09$,
both held fixed across every row. Adjoint-vector fields are stored
under the linear layout $v[\,12 x + 3\mu + a\,]$ on $12V$ entries,
$V = L_t L_s^3$, and the three operators entering the branch-response
determinant,
\begin{equation}
\Delta_1 \;=\; -D^2 + 2F + m^2\,\mathbb{1},
\qquad
\Delta_0 \;=\; -D^2 + m^2\,\mathbb{1},
\qquad
K_0 \;=\; -D_0^2,
\label{eq:operators-summary}
\end{equation}
are implemented with the sign and regulator conventions fixed in
Appendix~\ref{app:conventions}. Infrastructure self-tests verifying
consistency of the code paths entering~\eqref{eq:operators-summary}
are summarized in Appendix~\ref{app:infra-tests}. The computation is
carried out under the parity-even hypotheses stated in
Section~\ref{sec:phase}. No further sectoral structure enters the
numerical pipeline below.



\section{Physical Jacobian}
\label{sec:jacobian}

\subsection{Left-trivialized tangents and the physical 5-manifold}
\label{subsec:jacobian-LT}

The HS caloron of Section~\ref{sec:lattice-setup} is characterized by
five collective coordinates: four translations $z^\mu$ ($\mu=0,1,2,3$)
and one scale $\rho$. The corresponding tangents on the link manifold
are constructed in the left-trivialized frame, in which the variation
of a link $U_{x,\mu}$ is written as an element of the Lie algebra
$\mathfrak{su}(2)$ at the origin of the group via
$\delta U_{x,\mu} = U_{x,\mu}\,t_{x,\mu}$, so that
$t = \{t_{x,\mu}^a\}$ lives in the same adjoint-vector space
$\mathbb{C}^{12V}$ as the operators of
Section~\ref{sec:lattice-setup}. Each collective-coordinate tangent is
then a finite-difference derivative of the HS seed along the
corresponding direction:
\begin{equation}
t_\mu^{(\text{tr})}
\;\equiv\;
\frac{\mathrm{LT}\bigl[U(z+\delta\hat{e}_\mu)\bigr]
      - \mathrm{LT}\bigl[U(z-\delta\hat{e}_\mu)\bigr]}
     {2\,\delta},
\qquad
t^{(\text{sc})}
\;\equiv\;
\frac{\mathrm{LT}\bigl[U(\rho+\delta)\bigr]
      - \mathrm{LT}\bigl[U(\rho-\delta)\bigr]}
     {2\,\delta},
\label{eq:LT-tangents}
\end{equation}
with $\mathrm{LT}[\,\cdot\,]$ the left-trivialization map defined
in~\cite{companion-structural} and finite-difference steps
$\delta = 0.1$ for translations and $\delta = 0.05$ for scale. The
five-dimensional subspace
$V_{\rm mod} = \mathrm{span}\{\,t_0^{(\text{tr})},t_1^{(\text{tr})},
t_2^{(\text{tr})},t_3^{(\text{tr})},t^{(\text{sc})}\,\}$ is the
physical moduli-space tangent subspace at the caloron configuration.
Its implementation in \texttt{jacobian\_left\_triv.py} is exercised by
the infrastructure tests of Appendix~\ref{app:infra-tests}.

\subsection{Gauge cleaning and orientation collapse}
\label{subsec:jacobian-gauge-cleaning}

On a periodic $T^4$ with trivial holonomy, the three colour-orientation
generators $[T^a,\,A_\nu^{\rm inst}]$ produce tangents that are pure
gauge: they lie inside the gauge-orbit tangent space at the caloron
configuration and carry no physical variation of the collective
coordinates. The raw left-trivialized tangents
of~\eqref{eq:LT-tangents} contain small residual components along
those gauge directions, arising from discretization asymmetries of
the finite-difference step. They are removed by the orthogonal
projector
$P_\perp = \mathbb{1} - V_{\rm gauge}V_{\rm gauge}^{\!\top}$ defined
in Appendix~\ref{app:conventions}, with $V_{\rm gauge}$ an
orthonormal basis of the gauge-orbit tangent space, implemented in
\texttt{jacobian\_exact.py:gauge\_clean\_tangent}. After cleaning, the
residual gauge overlap of the physical tangents, measured by the
orientation indicator
\begin{equation}
G_{\rm orient}
\;\equiv\;
\max_{a=1,2,3}\;\max_{t\,\in\,V_{\rm mod}}\,
\bigl|\langle\,[T^a,A_\nu^{\rm inst}]\,|\,t\rangle\bigr|^2,
\label{eq:Gorient}
\end{equation}
satisfies $G_{\rm orient} \simeq 4\times 10^{-6}$ on every production
background. The eight raw tangents of the full $4 + 1 + 3$
decomposition therefore collapse to the five physical modes of
Section~\ref{subsec:jacobian-LT}, consistent with the fact that on
$T^4$ the three orientation directions carry no physical moduli.

\subsection{The $5\times 5$ Gram and the measured Jacobian}
\label{subsec:jacobian-Gram}

The physical Jacobian of the collective-coordinate measure is the
volume-form normalization of the five cleaned tangents, evaluated as
the positive square root of the determinant of the $5\times 5$ Gram
matrix
\begin{equation}
G_{ab}
\;=\;
\bigl\langle\,P_\perp\,t_a\;|\;P_\perp\,t_b\,\bigr\rangle,
\qquad
a,b \in \bigl\{z^0,z^1,z^2,z^3,\rho\bigr\},
\label{eq:Gram-5x5}
\end{equation}
with the inner product the standard $\mathbb{C}^{12V}$ inner product.
The matrix $G_{ab}$ is positive-definite (test~9 of
Appendix~\ref{app:infra-tests}), and the Jacobian is
\begin{equation}
\boxed{\;\Jphys \;\equiv\; \sqrt{\det G_5} \;=\; +\,0.118164\;}
\label{eq:Jphys-value}
\end{equation}
on Row~A, with Row~B within the volume-stability systematic of
Section~\ref{subsec:jacobian-systematics}. The sign of $\Jphys$ is
fixed by the positive square root of a positive-definite Gram
determinant. All subsequent uses of $\Jphys$ in the paper refer to
the value~\eqref{eq:Jphys-value}; it is not re-fitted or re-normalized
elsewhere.

\subsection{Systematics}
\label{subsec:jacobian-systematics}

Five systematic effects have been audited on the production rows.
Table~\ref{tab:jacobian-systematics} records the observed deviation of
$\Jphys$ under each, measured on Row~A and confirmed within the
stated tolerance on Row~B.

\begin{table}[h]
\centering
\begin{tabular}{lll}
\toprule
Effect & Diagnostic & Size \\
\midrule
Volume stability
  & $|\Jphys(\text{Row A}) - \Jphys(\text{Row B})|/\Jphys$
  & $0.4\%$ \\
Finite-difference step
  & convergence $\delta \to \delta/2$
  & $0.007\%$ \\
Gauge-cleaning residual
  & cleaning ratio
    $\langle t|V_{\rm gauge}V_{\rm gauge}^{\!\top}|t\rangle
     / \langle t|t\rangle$
  & $< 10^{-6}$ \\
Gram conditioning
  & $\kappa(G_5) = \lambda_{\max}/\lambda_{\min}$
  & $1.27$ \\
$F_{\mu\nu}$ correlation
  & normalized seeded--clover field-strength overlap
  & $0.92$--$0.94$ \\
\bottomrule
\end{tabular}
\caption{Jacobian systematics on Row~A. Values are representative of
Row~B within the volume-stability deviation listed in the first row.}
\label{tab:jacobian-systematics}
\end{table}

The largest deviation, the $0.4\%$ variation of $\Jphys$ between
Rows~A and~B, is subdominant to the statistical uncertainty on the
branch-response bracket reported in
Section~\ref{sec:finite-xi}. It is therefore retained only as a
sub-dominant contribution to the assembled master-formula uncertainty
of Section~\ref{sec:assembled}.



\section{Phase theorem}
\label{sec:phase}

\subsection{Parity-even class}
\label{subsec:phase-class}

The phase factor $\cos\Theta_\gamma$ in the master
formula~\eqref{eq:master} is fixed by the discrete symmetry content
of the microscopic sector whose one-loop integration generates the
gravitational effective action. We state the five hypotheses under
which this phase is determined analytically. These hypotheses are
framework-level conditions established
in~\cite{companion-structural} and are assumed throughout the
numerical pipeline of this paper.

\begin{enumerate}
\item[(H1)] \emph{Parity-even proto-field sector.} The microscopic
proto-field action is invariant under the spatial parity involution
$P : (x^0,\vec{x})\mapsto(x^0,-\vec{x})$.
\item[(H2)] \emph{Real-scalar labels.} All proto-field components are
real scalars; no complex or pseudoscalar component enters the
fundamental action.
\item[(H3)] \emph{Parity-even coupling to the emergent gauge field.}
The interaction of the proto-field sector with the emergent
coframe--gauge field is parity-even and preserves the involution
of (H1).
\item[(H4)] \emph{Parity-preserving lattice implementation.} The
lattice discretization and infrared regulator used throughout the
pipeline (symmetric clover field strength, Hermitian operators,
infrared regulator $m^2 = 0.01$) are invariant under~$P$.
\item[(H5)] \emph{No bare $\theta$-term.} The microscopic action
contains no bare topological-charge coupling $\theta\,Q$.
\end{enumerate}

Hypotheses (H1)--(H5) define the parity-even class in which the
phase theorem below applies. In the present paper these hypotheses
are not re-derived; they define the class of microscopic sectors for
which the lattice computation is carried out.

\begin{theorem}[Parity-even phase]
\label{thm:phase}
Let the microscopic proto-field sector be parity-even, real-scalar,
coupled to the emergent gauge field through a parity-even
interaction, with a parity-preserving regulator and no bare
$\theta$-term. Then the induced $\theta$-angle vanishes identically,
\[
\theta_{\rm ind}(X) \;=\; 0 \qquad \forall X,
\]
and the phase of the instanton contribution to the master formula is
\[
\cos\Theta_\gamma \;=\; +1
\]
on $\bR^4$, $S^1\times\bR^3$, and periodic tori $T^4$. No
simple-connectedness or dilute-instanton-gas assumption is required.
\end{theorem}

\begin{proof}
Under the spatial parity involution
$P: (x^0,\vec{x}) \mapsto (x^0,-\vec{x})$, the CP-odd density
$F_{\mu\nu}\tilde{F}^{\mu\nu}$ is odd:
$P^*[F\tilde F] = -F\tilde F$. Hypotheses (H1)--(H4) imply that the microscopic action, the
functional measure, and the regulator are each invariant under~$P$,
so the induced effective action $\Gamma_{\rm ind}[A]$ obtained by
integrating out the proto-field sector is a parity-even functional
of the gauge background~$A$. A
parity-even functional admits no term proportional to the parity-odd
density $F\tilde F$; consequently the induced $\theta$-angle, which
is the coefficient of $\int F\tilde F$ in $\Gamma_{\rm ind}[A]$,
vanishes identically:
\[
\theta_{\rm ind}(X) \;=\; 0 \qquad \text{for all } X.
\]
Hypothesis (H5) ensures no bare $\theta$-term is present to add to
the induced one. The phase of the instanton contribution to the
vacuum functional is $e^{i\theta_{\rm ind} Q}$ on a charge-$Q$
background, so $\cos\Theta_\gamma = +1$ on every admissible
background, independent of the base-manifold topology. The argument
uses only the parity transformation of the CP-odd operator
$F\tilde F$ and the parity invariance of the proto-field action; no
simple-connectedness or dilute-instanton-gas assumption is invoked.
\end{proof}

\subsection{Lattice consistency check}
\label{subsec:phase-consistency}

The theorem is consistent with direct lattice measurement: the
imaginary part of the branch-response trace $\delta_\xi \log\zeta$
on every production row of Section~\ref{sec:finite-xi} is below the
Hutchinson noise floor,
$|\mathrm{Im}(\delta_\xi \log\zeta)| < 10^{-8}$, so that the
production values are numerically real at a level far below the
quoted statistical uncertainty, and no empirical correction to
$\cos\Theta_\gamma = +1$ is required.


\section{Spectral structure and adaptive-prime prescription}
\label{sec:spectral}

\subsection{Three-tier decomposition}
\label{subsec:three-tier}

On the three production backgrounds of
Section~\ref{sec:lattice-setup}, the low spectrum of the lattice
vector operator $\Delta_1$ separates empirically into three tiers,
distinguished by the joint localisation--eigenvalue coordinates
$(\mathrm{IPR}\!\times\!V,\,\lambda)$, where
\begin{equation}
\mathrm{IPR}(\psi) \;\equiv\;
\frac{\sum_{x,\mu,a}|\psi_\mu^a(x)|^4}
     {\bigl[\sum_{x,\mu,a}|\psi_\mu^a(x)|^2\bigr]^2}
\label{eq:IPR-def}
\end{equation}
is the inverse-participation ratio of an adjoint-vector eigenmode
and $V = L_t L_s^3$ is the lattice volume. Diagonalization of the
$40$ lowest-lying eigenmodes of $\Delta_1$ identifies the three
tiers directly in the $(\mathrm{IPR}\!\times\!V,\,\lambda)$ plane
(Figure~\ref{fig:three-tier}).
\begin{itemize}
\item \emph{Tier 1 (torus).}
Twelve delocalised modes with $\mathrm{IPR}\!\times\!V$ in the range
$[1.0,\,1.08]$ and eigenvalues $\lambda \in [0.006,\,0.018]$.
\item \emph{Tier 2 (core).}
Eight localised modes with $\mathrm{IPR}\!\times\!V$ in the range
$[14,\,161]$ and eigenvalues $\lambda \in [0.016,\,0.030]$.
\item \emph{Tier 3 (bulk).}
All modes with $n \geq 20$,
$\mathrm{IPR}\!\times\!V \in [1.2,\,1.6]$, and eigenvalues
separated from Tier~2 by a clean gap
$\delta_{\rm bulk} \in [0.049,\,0.084]$, a factor of five to eight
times the infrared regulator $m^2$.
\end{itemize}

\begin{figure}[h]
\centering
\includegraphics[width=0.65\textwidth]{fig_three_tier_rowA.pdf}
\caption{Joint distribution of $(\mathrm{IPR}\!\times\!V,\,\lambda)$
for the $40$ lowest-lying eigenmodes of $\Delta_1$ on Row~A. The
three tiers separate cleanly in both axes: Tier~1 (torus, $12$
delocalised modes) occupies $\lambda \in [0.006,\,0.018]$ with
$\mathrm{IPR}\!\times\!V \in [0.12,\,0.51]$; Tier~2 (core, $8$
localised modes) occupies $\lambda \in [0.028,\,0.030]$ and further
splits by $\mathrm{IPR}\!\times\!V$ into a low-subset of four modes
($\mathrm{IPR}\!\times\!V \simeq 2.4$) identified with the translation
quartet and a high-subset of four modes
($\mathrm{IPR}\!\times\!V \in [18.2,\,23.8]$) identified with the
scale singlet and gauge triplet; Tier~3 (bulk) sits above the bulk
gap at $\lambda \simeq 0.12$. The dotted line marks the infrared
regulator $m^2 = 0.01$.}
\label{fig:three-tier}
\end{figure}

\begin{table}[h]
\centering
\begin{tabular}{lcccc}
\toprule
Row & torus $\mathrm{IPR}\!\times\!V$ &
\multicolumn{2}{c}{core $\mathrm{IPR}\!\times\!V$} &
bulk gap $\delta_{\rm bulk}$ \\
\cmidrule(lr){3-4}
 & & low subset (4) & high subset (4) & \\
\midrule
A & $0.12$--$0.51$ & $2.42$--$2.47$ & $18.2$--$23.8$ & $0.084$ \\
B & $0.12$--$0.50$ & $3.45$--$3.49$ & $25.6$--$35.6$ & $0.050$ \\
C & $0.13$--$0.50$ & $4.26$--$4.29$ & $31.8$--$41.4$ & $0.031$ \\
\bottomrule
\end{tabular}
\caption{Measured localisation and bulk-gap structure on the three
production rows. The core tier splits empirically into two subsets of
four modes each; the low-$\mathrm{IPR}\!\times\!V$ subset is identified
with the translation quartet and the high-$\mathrm{IPR}\!\times\!V$ subset
with the scale singlet and gauge triplet (Section~\ref{subsec:V20-content}).
The normalization is $\mathrm{IPR}\!\times\!V$ with $V = L_t L_s^3$;
values are read directly from the production pipeline
\texttt{mode\_diagnostic.py}. The row-to-row growth of the high-subset
$\mathrm{IPR}\!\times\!V$ reflects the sharper localisation of scale and
gauge modes on finer lattices.}
\label{tab:three-tier-summary}
\end{table}

On each of the three production rows, Tiers~1 and~2 are separated by
a clean eigenvalue gap (Table~\ref{tab:three-tier-summary}, columns
2--4) and by a further separation in $\mathrm{IPR}\!\times\!V$
(localisation). While the eigenvalue-level separation is quantitative,
it is narrow: the Tier~1/Tier~2 gap narrows from $\lambda_2^{\rm min}
- \lambda_1^{\rm max} \simeq 0.010$ on Row~A to $\simeq 0.002$ on
Row~C. Naive eigenvalue-only separation criteria are therefore
fragile on larger lattices; the IPR-based identification of
§\ref{subsec:adaptive-prime} is robust across the full production
range. Across all three rows the torus/core cardinality $(12,\,8)$,
the IPR separation into a low-subset ($4$ translation-quartet modes)
and a high-subset ($4$ scale+gauge modes), and the existence of a
clean bulk gap (decreasing from $\delta_{\rm bulk} = 0.084$ on Row~A
to $0.031$ on Row~C) are all preserved.

\subsection{Content of the low manifold}
\label{subsec:V20-content}

The twenty-dimensional low manifold
\[
V_{20} \;\equiv\; V_{\rm tor}^{12} \,\oplus\, V_{\rm core}^{8}
\]
splits numerically into a torus sector and a core sector. This
splitting is checked by singular-value decomposition of
$(\mathbb{1} - P_{\rm tor})V_{20}$, whose first eight singular values
are numerically indistinguishable from unity on every production
row.

\paragraph{Torus sector ($V_{\rm tor}^{12}$).}
The twelve torus modes carry the degeneracy pattern $(3+5+1+3)$
predicted by the $J_1 \!\cdot\! J_2$ Casimir decomposition of the
self-dual Weitzenb\"ock endomorphism acting on the constant-one-form
subspace~\cite{companion-structural}. Direct measurement on every
production row reproduces the expected multiplicities $(3, 5, 1, 3)$
on $V_{\rm tor}^{12}$, with a numerically traceless restricted
curvature spectrum and the predicted two-cluster Casimir pattern.
Projection onto the free constant-one-form subspace captures
$97$--$98\%$ of the torus norm, similarly across the three
production rows.

\paragraph{Core sector ($V_{\rm core}^{8}$).}
The eight core modes are localised near the caloron core and
further decompose into a five-dimensional physical moduli-space
block and a three-dimensional orientation block:
\[
V_{\rm core}^{8}
\;=\;
V_{\rm mod}^{5}
\,\oplus\,
V_{\rm orient}^{3}.
\]
The moduli-space block comprises the four translation modes and the
one scale mode of Section~\ref{sec:jacobian}; the orientation block
is spanned by the three longitudinal constant-colour-rotation
directions $[T^a,\,A_\nu^{\rm inst}]$. The three orientation modes
are the localised core counterparts of the constant colour-rotation
directions cleaned in Section~\ref{subsec:jacobian-gauge-cleaning}. They
are the modes associated with the longitudinal sector that are
removed in the gauge-invariant one-loop determinant by the
combination $-\tfrac{1}{2}\Tr_{\rm vec} + \Tr_{\rm gh}$.

\paragraph{Identification via production overlaps.}
Sector membership within $V_{\rm core}^{8}$ is identified by the
overlaps with the production tangent bases. On every admissible
background the measured pattern is as follows: torus modes have
$O_{\rm tor} \simeq 0.97$ and vanishing overlap with the
moduli-space and orientation bases; the scale singlet has
$O_{\rm LT} \simeq 0.92$; the translation quartet has
$O_{\rm LT} \in [0.22,\,0.26]$; and the orientation triplet has
$O_{\rm gauge} \in [0.54,\,0.55]$. The translation-quartet and
orientation-triplet undercaptures are reproducible across the three
production rows. They are a structural feature of the caloron
zero-mode dressing on the compact torus, not a numerical artifact.

\subsection{Adaptive-prime construction}
\label{subsec:adaptive-prime}

The branch-response determinant of Section~\ref{sec:finite-xi} is
evaluated on the orthogonal complement of the low manifold. Let
$V_{20}$ be the $12V \!\times\! 20$ matrix whose columns are the
orthonormal low-manifold basis vectors identified in
Section~\ref{subsec:three-tier} by the IPR criterion. The adaptive-prime
projector is
\begin{equation}
P \;=\; \mathbb{1} \,-\, V_{20}\,V_{20}^{\!\top},
\qquad
P^2 \;=\; P,\qquad
P\,V_{20} \;=\; 0,
\label{eq:adaptive-prime-projector}
\end{equation}
and the primed trace entering the first-order integrand $h(\xi)$ is
\begin{equation}
\Tr'\!\bigl[\Delta_1^{-1}\,A'\bigr]
\;\equiv\;
\Tr\!\Bigl[P\,(P\,\Delta_1\,P)^{-1}\,P\,A'\,P\Bigr],
\label{eq:primed-trace}
\end{equation}
computed by projected conjugate gradient on the complement of
$V_{20}$. The $A'$ insertion is the anisotropic operator derivative
of Appendix~\ref{app:conventions}. At each finite-$\xi$ Simpson
node, $V_{20}(\xi)$ is obtained from $V_{20}(1)$ by the transported
construction of Appendix~\ref{app:transported-prime}; the projector
$P(\xi)$ inherits the same transport.

\subsection{Resolvent-Ritz dressing of moduli tangents}
\label{subsec:resolvent-Ritz}

The undercaptures reported in Section~\ref{subsec:V20-content} reflect the
fact that bare continuum moduli tangents do not coincide with the
true low eigenmodes of $\Delta_1$ on a compact torus. In particular,
the translation quartet is substantially dressed by compact-geometry
effects and is therefore only partially captured by bare
finite-difference tangents. Numerical tests show that a
resolvent-Ritz refinement of the bare tangent basis can recover the
same low-manifold directions to better than $99.9\%$ overlap on the
production rows. Because this refinement is not part of the
production determinant pipeline, and because the present paper does
not require it operationally, its full implementation is deferred to
separate technical work.

\subsection{Validation at the level of the branch response}
\label{subsec:operational-closure}

For the numerical purposes of the present paper, the adaptive-prime
prescription is validated at the level of the branch response
$\delta_\xi \log\zeta$ by three independent checks, each reported
later in the paper:
\begin{enumerate}
\item[(i)] lattice-size stability across the three production rows
spanning a factor of $11$ in lattice volume
(Section~\ref{sec:continuum});
\item[(ii)] cross-row agreement at the sub-$0.2\sigma$ level on
every pair (Section~\ref{sec:finite-xi});
\item[(iii)] independent-scheme agreement with the
projected-complement exact prime of Section~\ref{subsec:three-prescriptions}
at $0.77\sigma$ on Row~A with $800$ paired samples
(Section~\ref{subsec:cross-check-exact}).
\end{enumerate}
The combination of the three checks is stronger than any one
individually, because they test different failure modes: (i)
finite-spacing artefacts, (ii) statistical-systematic mismatch
between rows, and (iii) prescription-dependence of the final
bracket.

The remaining open technical item is a first-principles dressed
analytic basis: an explicit representation-theoretic construction
of the compact-$T^4$ dressing of the moduli-space tangents that
would replace the IPR-based identification of $V_{20}$ by a
closed-form basis. This is not required for the numerical result
reported here. It is a separate technical refinement of the
low-manifold construction.

\section{Prescription history and cross-validation}
\label{sec:history}

The primed trace~\eqref{eq:primed-trace} depends on the choice of
low-mode subspace removed from $\Delta_1^{-1}$. In this section we
name three prescriptions that are in principle applicable to the
branch-response determinant~\eqref{eq:master}, state which is used
for the production result, and report an independent cross-check of
the production prescription against one of the alternatives.

\subsection{The three prescriptions}
\label{subsec:three-prescriptions}

\paragraph{Threshold prime.}
The simplest prescription removes all modes of $\Delta_1$ below a
hard eigenvalue threshold $\lambda_*$ on the instanton background
and all exact zero modes on the free background. Its variance is
the smallest of the three prescriptions considered here, at the
cost of a single empirical scale $\lambda_*$ whose placement is
not unambiguous on backgrounds where Tiers~1 and~2 overlap in
eigenvalue (Section~\ref{subsec:three-tier}). Because the present paper
values prescription clarity over raw precision, the threshold
prime is not used here; we note it only for completeness.

\paragraph{Projected-complement exact prime.}
A second prescription uses only the five physical moduli-space
directions as the removed subspace: the gauge-cleaned
left-trivialized translation and scale tangents of
Section~\ref{sec:jacobian}. Let $V_{\rm mod}$ denote the
$12V \times 5$ matrix whose columns are an orthonormal basis
obtained from those five cleaned tangents. Then the projector onto
the complement of the physical moduli subspace is
\begin{equation}
P_{\rm mod} \;=\; \mathbb{1} \,-\, V_{\rm mod}\,V_{\rm mod}^{\!\top},
\label{eq:Pmod}
\end{equation}
and the constrained Green operator is
\begin{equation}
G_{P_{\rm mod}}
\;=\;
P_{\rm mod}\,\bigl(P_{\rm mod}\,\Delta_1\,P_{\rm mod}\bigr)^{-1}
P_{\rm mod}.
\label{eq:exact-prime-Green}
\end{equation}
This construction uses no IPR input, no threshold, and no
low-manifold identification beyond the physical collective
coordinates; it is the prescription with the fewest empirical
inputs. Its limitation is the Hutchinson variance inherited from
leaving the twelve torus and three orientation modes of
Section~\ref{subsec:V20-content} in the trace. We use it below as a direct
cross-check of the production prescription.

\paragraph{Adaptive transported prime.}
The production prescription is the adaptive-prime construction of
Section~\ref{subsec:adaptive-prime}, transported across the Simpson grid
by the algorithm of Appendix~\ref{app:transported-prime}. It has
smaller Hutchinson variance than~\eqref{eq:exact-prime-Green}
because the removed subspace additionally contains the twelve torus
and three orientation modes identified in Section~\ref{subsec:V20-content};
at the same precision, fewer noisy low-eigenvalue excitations
contribute to the trace. The numerical values reported in
Section~\ref{sec:finite-xi} use this prescription.

\subsection{Cross-validation: projected-complement exact prime}
\label{subsec:cross-check-exact}

The adaptive prime is identified by an
inverse-participation-ratio criterion rather than by a closed-form
analytic construction. A natural concern is
prescription-circularity: the consistency of the production result
across three rows (Sections~\ref{sec:finite-xi}
and~\ref{sec:continuum}) could in principle be a property of the
adaptive prescription rather than of the underlying branch
response. We close this concern by comparing the adaptive prime to
the projected-complement exact prime
of~\eqref{eq:Pmod}--\eqref{eq:exact-prime-Green}, which shares no
empirical inputs with the adaptive construction beyond the five
physical moduli tangents.

On Row~A with $N_S = 800$ paired Hutchinson samples at the common
seed of Appendix~\ref{app:conventions}, the two prescriptions yield
\begin{equation}
\delta_\xi \log\zeta\bigr|_{\rm exact\ prime}^{\rm Row\,A}
\;=\; -2.49 \pm 1.06,
\qquad
\delta_\xi \log\zeta\bigr|_{\rm adaptive\ prime}^{\rm Row\,A}
\;=\; -1.66 \pm 0.19,
\label{eq:scheme-comparison}
\end{equation}
with scheme difference $|\Delta| = 0.83$, combined statistical
uncertainty $1.07$, and tension $0.77\sigma$. 
The adaptive-prime value in~\eqref{eq:scheme-comparison} is the
Row~A entry of Table~\ref{tab:production-results} at $N_S = 80$; the
exact-prime value is the result of a dedicated $N_S = 800$ campaign
on the same row, with the common seed of
Appendix~\ref{app:conventions}.
The exact-prime
result corresponds to the bracket $\Bobs^{\rm exact} = 0.917$,
while the adaptive-prime result gives
$\Bobs^{\rm adapt} = 0.810$; the difference is discussed below only
as a consequence of the nonlinear map
$\Bobs = 1 - e^{\delta_\xi \log\zeta}$, not as an inconsistency of
the underlying $\delta_\xi \log\zeta$ values. The two prescriptions
agree on sign, on order of magnitude, and on statistical
consistency within their combined error.

The larger uncertainty on the exact-prime result is a structural
consequence of the prescription, not a numerical limitation.
Leaving the twelve torus and three orientation modes in the trace
preserves the lowest-eigenvalue excitations of $\Delta_1$ on the
caloron, and through the inverse operator these dominate the
Hutchinson-estimator variance. The adaptive prime achieves a
tighter error precisely because it additionally removes those
fifteen modes, at the price of an empirically identified low
manifold. The two prescriptions thus sit at opposite ends of a
clarity-of-construction--versus--precision trade-off, and their
agreement within the exact prime's statistical error is the
strongest prescription-independent check available at the present
precision.

\begin{remark}[Bracket representation]
\label{rem:bracket-nonlinearity}
The scheme difference $0.83$ in $\delta_\xi \log\zeta$ translates
via $\Bobs = 1 - e^{\delta_\xi \log\zeta}$ into a bracket
difference
$\Bobs^{\rm exact} - \Bobs^{\rm adapt} = 0.917 - 0.810 = 0.107$.
The nonlinearity of the bracket map amplifies the $\delta$-level
agreement into a larger apparent bracket offset. At current
precision, the offset sits within the combined statistical
uncertainty on $\delta_\xi \log\zeta$ and should be read at the
$\delta$ level rather than at the bracket level.
\end{remark}

The three independent checks of Section~\ref{subsec:operational-closure}
accordingly read as follows in quantitative form. Cross-row
agreement (Section~\ref{sec:finite-xi}) is sub-$0.2\sigma$ on every pair;
lattice-size stability (Section~\ref{sec:continuum}) holds across a
factor of $11$ in volume; and the cross-prescription agreement
measured here is $0.77\sigma$ on Row~A with $800$ paired samples.
Taken together, the three checks provide the strongest available
validation of the production prescription; no one of them would be
as convincing in isolation. Of the three, the universality
statement of Proposition~\ref{prop:universality} carries the
strongest empirical weight, because it is measured against noise
at the percent level across an order of magnitude in lattice
volume, whereas the scheme check reported here is limited by the
structural variance of removing only the physical moduli.


\section{Finite-$\xi$ campaign: transported-prime Simpson}
\label{sec:finite-xi}

\subsection{Simpson scheme and transported prime}
\label{subsec:simpson-setup}

The branch-response determinant of Section~\ref{sec:spectral} is
required at $\xi \in [1,\,1+\epsilon]$ with $\epsilon = 0.09$. In the
production convention, the assembled response is the sum of a Simpson
quadrature of the $\xi$-dependent integrand and a $\xi$-independent
anchor-Jacobian contribution:
\begin{equation}
\delta_\xi \log\zeta
\;=\;
\frac{\epsilon}{6}\,\Bigl[\,h(1) + 4\,h(1+\epsilon/2) + h(1+\epsilon)\,\Bigr]
\;+\; J_{\rm anchor},
\label{eq:simpson-quadrature}
\end{equation}
where $h(\xi)$ is the Jacobian-free branch-response integrand at the
indicated $\xi$ and $J_{\rm anchor}$ is the log-Jacobian contribution
of the left-trivialized collective-coordinate measure evaluated on the
caloron background. In the free-background subtraction convention of
Appendix~\ref{app:conventions}, the anchor-Jacobian contribution is
$\xi$-independent and coincides numerically with the physical
Jacobian of Section~\ref{sec:jacobian},
\begin{equation}
J_{\rm anchor} \;=\; \Jphys \;=\; 0.118164,
\label{eq:Janchor-equals-Jphys}
\end{equation}
on every production row. With this convention, the first-order
estimator is $\delta_\xi \log\zeta|_{\rm 1st} = \epsilon\,h(1) +
J_{\rm anchor}$, which reproduces the production values in
Table~\ref{tab:production-results}.

The equality~\eqref{eq:Janchor-equals-Jphys} is a convention choice,
not a coincidence: the left-trivialized collective-coordinate measure
enters the partition function multiplicatively through a
Gram-determinant factor $|\det G|^{1/2} = \Jphys$, and the
free-background subtraction used in Appendix~\ref{app:conventions}
places the $\xi$-independent part of the associated log-Jacobian
contribution on the right-hand side of the Simpson
formula~\eqref{eq:simpson-quadrature} rather than inside the
integrand $h(\xi)$. Alternative subtraction conventions redistribute
this contribution between $h$ and the additive constant; under the
convention fixed here, the additive constant coincides numerically
with $\Jphys$, so that the same quantity appears in both
Section~\ref{sec:jacobian} and the Simpson assembly. The assembled
$\delta_\xi \log\zeta$ is convention-independent.

At each of the two non-anchor Simpson nodes
$\xi \in \{1.045,\,1.090\}$, the adaptive-prime subspace
$V_{20}(\xi)$ is obtained from $V_{20}(1)$ by the transported
construction of Appendix~\ref{app:transported-prime}. The
construction re-projects onto the eigenbasis of $\Delta_1(\xi)$ at
the transported node and assigns mode identities by greedy overlap
matching against the anchor, so that the twenty low modes tracked
at each Simpson node correspond by continuity to the twenty anchor
modes at $\xi = 1$. The transported projector
$P(\xi) = \mathbb{1} - V_{20}(\xi)V_{20}(\xi)^{\!\top}$
of Section~\ref{subsec:adaptive-prime} inherits the same transport,
so that the primed trace $\Tr'[\Delta_1^{-1}(\xi)\,A'(\xi)]$ is
defined consistently across the Simpson grid.

Each node integrand $h(\xi_j)$ is computed as a Hutchinson trace
with $N_S = 80$ paired Rademacher vectors at the common noise seed
of Appendix~\ref{app:conventions}. The paired construction shares
the same noise realisation between the anchor and each transported
node, so that the variance of the Simpson estimator is dominated by
the variances of the paired differences $h(\xi_j) - h(\xi_0)$
rather than by the variances of the absolute node values. Linear
solves against the projected operator
$P(\xi)\,\Delta_1(\xi)\,P(\xi)$ are performed on an NVIDIA H100
$80$\,GB GPU by \texttt{cupyx.scipy.sparse.linalg.cg} wrapping the
CuPy \texttt{LinearOperator} of Appendix~\ref{app:conventions}, to
a relative-residual tolerance of $10^{-10}$. Per-row wall time on
the GPU pipeline is approximately $25$ minutes at the production
configuration.

\subsection{Production results}
\label{subsec:production-results}

\begin{table}[h]
\centering
\begin{tabular}{lcccccc}
\toprule
Row & $\delta_\xi \log\zeta$ (1st) & $\delta_\xi \log\zeta$ (Simpson) &
err & $\Delta_{\rm NL}$ & $\Bobs$ & $\Icl$ \\
\midrule
A: $18^3\times 12$, $\rho=3.0$
  & $-1.928$ & $-1.660$ & $0.186$ & $+13.9\%$ & $0.810$ & $0.983$ \\
B: $24^3\times 16$, $\rho=3.5$
  & $-1.959$ & $-1.695$ & $0.213$ & $+13.5\%$ & $0.816$ & $0.998$ \\
C: $32^3\times 24$, $\rho=4.5$
  & $-1.969$ & $-1.646$ & $0.275$ & $+16.4\%$ & $0.807$ & $1.032$ \\
\bottomrule
\end{tabular}
\caption{Transported-prime Simpson results on the three production
rows. The three rows span a factor of $11\times$ in lattice volume
and a factor of $1.5$ in caloron size.}
\label{tab:production-results}
\end{table}

\paragraph{Pairwise consistency.}
All three pairwise Simpson tensions are below $0.2\sigma$: A versus
B at $0.12\sigma$, A versus C at $0.04\sigma$, B versus C at
$0.14\sigma$ (see Table~\ref{tab:pairwise} in
Section~\ref{sec:continuum}).

\paragraph{Weighted result.}
Inverse-variance weighted across the three rows,
\begin{equation}
\delta_\xi \log\zeta \;=\; -1.669 \pm 0.125,
\qquad
\Bobs \;=\; 0.811 \pm 0.023.
\label{eq:bracket-result}
\end{equation}
This is the bracket value that enters the assembled result of
Section~\ref{sec:assembled}.

\subsection{Anchor-branch normalization}
\label{subsec:zeta-normalization}

The master formula~\eqref{eq:master} contains the combination
$\Jphys\bigl[\zeta(0) - \zeta(X_*)\bigr]$. Writing this as
$\Jphys\,\zeta(0)\,\Bobs$ with
$\Bobs \equiv 1 - \zeta(X_*)/\zeta(0)$ measured in
\eqref{eq:bracket-result}, the assembly requires the anchor-branch
normalization $\zeta(0)$ evaluated on the caloron background. This
normalization is fixed by the production pipeline's convention:
$\zeta(0)$ is the value of the branch-response partition function
evaluated at $\xi = 1$ (the undeformed metric) on the caloron
background of Section~\ref{sec:lattice-setup}, in the free-background
subtraction convention of Appendix~\ref{app:conventions}. It is a
direct output of the same Hutchinson pipeline used for the Simpson
integrand, evaluated at $\xi = 1$ without any $\xi$-variation.

On Row~A, the production pipeline returns
\begin{equation}
\zeta(0) \;=\; 8.42 \pm 0.08,
\label{eq:zeta0-value}
\end{equation}
with the quoted $\pm 0.08$ the Hutchinson statistical uncertainty at
$N_S = 80$ on the anchor trace\footnote{The Row~A value of $\zeta(0)$ is the mean of the anchor
Hutchinson trace $\Tr'[\Delta_1^{-1}(\xi=1)\,A']$ over the
$N_S = 80$ paired samples, with the same common seed and
projected-CG tolerance as the Simpson integrand; it is read
directly from the production pipeline output and is not inferred
from the assembly arithmetic of Section~\ref{subsec:assembly-arithmetic}.}. The corresponding values on Rows~B and~C agree with~\eqref{eq:zeta0-value}
to within the per-row statistical uncertainty, consistent with the
sub-percent row-to-row variation of the Jacobian itself
(Section~\ref{subsec:jacobian-systematics}). Because $\zeta(0)$ is a
ratio-normalization factor in the master formula rather than an
independent physical observable, its volume independence at the
sub-percent level confirms that the anchor-branch normalization is
controlled at the level required for the assembly of
Section~\ref{sec:assembled}.

\subsection{Integrand values}
\label{subsec:integrand-values}

For referee reproducibility, Table~\ref{tab:simpson-integrands}
records the node-level integrand data from the dedicated
paired-noise campaign on Rows~A and~B, together with their paired
differences against the anchor $\xi_0 = 1$. Row~C enters the main
measurement through its final Simpson value in
Table~\ref{tab:production-results}; its node-level integrand values
are not used elsewhere in the paper. The integrands $h(\xi_j)$
decompose additively into the three contributions identified in
Section~\ref{sec:spectral},
\begin{equation}
h(\xi) \;=\; -\,\xi\,T(\xi)
\;-\; \tfrac{1}{2}\,F(\xi)
\;+\; 2\,\xi\,G(\xi),
\label{eq:integrand-decomposition}
\end{equation}
corresponding to the temporal-Laplacian, field-strength, and
ghost-operator contributions respectively. The paired differences
$h(\xi_j) - h(1)$ have tight errors an order of magnitude below
those of the absolute $h(\xi_j)$, consistent with the paired-noise
construction of Section~\ref{subsec:simpson-setup}.

\begin{table}[h]
\centering
\begin{tabular}{lccc}
\toprule
Row & $h(1.000)$ & $h(1.045) - h(1)$ & $h(1.090) - h(1)$ \\
\midrule
A: $18^3\times 12$, $\rho=3.0$
  & $-22.737 \pm 1.967$
  & $+3.090 \pm 0.127$
  & $+5.507 \pm 0.280$ \\
B: $24^3\times 16$, $\rho=3.5$
  & $-23.076 \pm 2.296$
  & $+3.149 \pm 0.192$
  & $+5.009 \pm 0.508$ \\
\bottomrule
\end{tabular}
\caption{Simpson-node integrand values on Rows~A and~B. Absolute
node values are reported only at the anchor $\xi_0 = 1$; the
non-anchor nodes are reported as paired differences against the
anchor, which have the tighter errors relevant for the Simpson
estimator. The first-interval slopes $[h(1.045) - h(1)]$ agree
across the two rows at $0.3\sigma$, independent evidence of the
cross-row stability established in
Table~\ref{tab:production-results}.}
\label{tab:simpson-integrands}
\end{table}

The row-to-row agreement of the paired differences is the more
relevant reproducibility diagnostic than the absolute anchor
values, because the paired differences are the quantities entering
the variance-reduced Simpson estimator.

\subsection{$\Icl$ trend and window}
\label{subsec:icl-correlation}

Rows outside the primary admissibility window
$\Icl \in [0.98,\,1.00]$ exhibit an empirical $\Icl$-dependence
reported in Appendix~\ref{app:icl-degraded}: a weighted linear
regression $\delta_\xi \log\zeta / \epsilon \approx -18.7 -
49.5\,(\Icl - 1)$ fitted to four reference rows with
$\Icl \in [1.07,\,1.33]$. Extrapolated to $\Icl = 1$, this gives
$\delta_\xi \log\zeta = -1.68 \pm 0.44$ (with the uncertainty
inflated by $\sqrt{\chi^2/\mathrm{dof}}$ to account for residual
scatter), consistent with the primary weighted result at the
$0.03\sigma$ level. The three production rows
of Table~\ref{tab:production-results} sit at
$\Icl \in [0.983,\,1.032]$, with Row~C at $\Icl = 1.032$ retained
as a stress test approximately $3\%$ outside the primary
admissibility window of Section~\ref{sec:lattice-setup}. Row~C's
pairwise tensions against Rows~A and~B ($0.04\sigma$ and
$0.14\sigma$ respectively, Table~\ref{tab:pairwise}) show that the
bracket value remains consistent even at this mild excursion from
the window. The primary admissibility window of
Section~\ref{sec:lattice-setup} is therefore retained as the formal
quality-control criterion; Row~C provides direct evidence that the
bracket is not fine-tuned to that window.
The linear guide of Appendix~\ref{app:icl-degraded} is fitted over
$\Icl \in [1.07,\,1.33]$ and does not extrapolate cleanly to
Row~C's $\Icl = 1.032$: naive linear extrapolation would predict
$\delta_\xi\log\zeta/\epsilon \approx -22.4$ at Row~C, whereas the
measurement gives $-18.3$. This mismatch is expected: the linear
form is descriptive within its fit range and loses predictive
power near $\Icl = 1$, where the caloron is no longer visibly
under-resolved. Row~C's consistency with Rows~A and~B at
sub-$0.2\sigma$ is therefore the stronger constraint, not the
trend's extrapolation.

\subsection{Stability of the nonlinear correction across three lattices}
\label{subsec:universality}

The Simpson estimator of $\delta_\xi \log\zeta$ differs from its
first-order Taylor approximation by a finite nonlinear correction
\begin{equation}
\Delta_{\rm NL} \;\equiv\;
\frac{\delta_\xi \log\zeta\bigr|_{\rm Simpson}
      - \delta_\xi \log\zeta\bigr|_{\rm 1st}}
     {\delta_\xi \log\zeta\bigr|_{\rm 1st}}.
\label{eq:Delta-NL}
\end{equation}
On the three production rows, $\Delta_{\rm NL}$ is measured as
\begin{align*}
\Delta_{\rm NL}(\text{Row A}) &= +13.9\%
  \qquad (\rho/L = 0.167,\ \Icl = 0.983),\\
\Delta_{\rm NL}(\text{Row B}) &= +13.5\%
  \qquad (\rho/L = 0.146,\ \Icl = 0.998),\\
\Delta_{\rm NL}(\text{Row C}) &= +16.4\%
  \qquad (\rho/L = 0.141,\ \Icl = 1.032),
\end{align*}
a three-row range of $2.9$ percentage points on a quantity of
unweighted mean $14.6\%$.

\begin{proposition}[Stability of the nonlinear correction]
\label{prop:universality}
Across three production rows spanning a factor of $11$ in lattice
volume, a factor of $1.5$ in caloron size, and $\Icl$ in the range
$[0.983, 1.032]$, the Simpson-to-first-order nonlinear correction
is measured in the range $\Delta_{\rm NL} \in [13.5\%,\,16.4\%]$
with unweighted mean $14.6\%$. The correction is independent of
$S_{\rm inst}$, of $\Jphys$, of $\cos\Theta_\gamma$, and of any
further sectoral input beyond the master formula: each of these
multiplies the master formula~\eqref{eq:master} and cancels in the
ratio~\eqref{eq:Delta-NL}. Within the primary admissibility window
$\Icl \in [0.98,\,1.00]$ the correction is stable at the
sub-percent level ($13.5\%$ vs $13.9\%$); the full three-row
spread of $2.9$ percentage points is larger than the statistical
error on any individual row but far smaller than the order-unity
drift a prescription-dependent artefact would produce across the
range of lattice volumes considered. This reading rests on the
interpretive assumption that a prescription-dependent artefact
would dominate the difference between lattices of substantially
different volumes, an assumption consistent with the scheme
cross-check of Section~\ref{subsec:cross-check-exact}, in which removing
a different low-mode subspace produces a $\delta$-level shift
contained within the structural variance of that prescription.
\end{proposition}

\begin{proof}
The quantity $\delta_\xi \log\zeta$ is computed directly from the
primed trace $\Tr'[\Delta_1^{-1}\,A']$ of
Eq.~\eqref{eq:primed-trace}, which depends only on the adjoint
vector operator $\Delta_1$ on the caloron background and on the
anisotropic insertion $A'$. The master-formula factors $S_{\rm
inst}$, $\Jphys$, and $\cos\Theta_\gamma$ multiply the assembled
cosmological constant but do not enter $\delta_\xi \log\zeta$ itself;
consequently $\Delta_{\rm NL}$, defined as a ratio of two
determinations of $\delta_\xi \log\zeta$ via
Eq.~\eqref{eq:Delta-NL}, is independent of these factors by
construction. What $\Delta_{\rm NL}$ does depend on is the
$\xi$-dependence of $\zeta$ itself, which is controlled by the
geometry of $\Delta_1$ on the caloron and by the smoothness of the
branch determinant as a function of $\xi$; these are properties
intrinsic to the branch-bias mechanism, not to its multiplicative
dressing. $\Delta_{\rm NL}$ is therefore the finite-$\xi$ observable
most directly diagnostic of the branch-bias mechanism that the
present pipeline provides.
\end{proof}

\begin{remark}[Why this matters]
\label{rem:universality-significance}
Had the nonlinear correction been dominated by a
prescription-dependent artefact --- for instance, an artefact of
the adaptive-prime projector or of the trace-estimator variance ---
one would expect a much larger drift between lattices differing by
an order of magnitude in volume. Its stability within a few
percentage points across such a range is therefore independent
evidence for the transported-prime construction, complementary to
the continuum consistency of Section~\ref{sec:continuum} and the
scheme cross-check of Section~\ref{subsec:cross-check-exact}.
\end{remark}

\subsection{Degraded rows}
\label{subsec:degraded-deferred}

Rows with $\Icl$ outside the range $[0.983,\,1.032]$ spanned by
Table~\ref{tab:production-results} are used only as a systematics
reference and are reported in Appendix~\ref{app:icl-degraded}; they
do not contribute to the weighted
bracket~\eqref{eq:bracket-result} or to any subsequent assembly.


\section{Continuum controls and lattice-size independence}
\label{sec:continuum}

\subsection{Three-lattice comparison}
\label{subsec:three-lattice}

Table~\ref{tab:production-results} reports three independent
Simpson measurements on lattices spanning lattice volumes
\begin{equation}
V_{\rm A} = 69\,120,
\qquad
V_{\rm B} = 221\,184,
\qquad
V_{\rm C} = 786\,432,
\label{eq:three-volumes}
\end{equation}
and caloron scale parameters $\rho \in \{3.0,\,3.5,\,4.5\}$. The
pairwise volume ratios are $V_{\rm B}/V_{\rm A} = 3.20$,
$V_{\rm C}/V_{\rm B} = 3.56$, and $V_{\rm C}/V_{\rm A} = 11.38$.
Table~\ref{tab:pairwise} summarises the pairwise Simpson tensions
across the three production rows.

\begin{table}[h]
\centering
\begin{tabular}{lccc}
\toprule
Pair & $|\Delta \delta_\xi \log\zeta|$ & combined $\sigma$ & tension \\
\midrule
A vs B & $0.035$ & $0.283$ & $0.12\sigma$ \\
A vs C & $0.014$ & $0.332$ & $0.04\sigma$ \\
B vs C & $0.049$ & $0.348$ & $0.14\sigma$ \\
\bottomrule
\end{tabular}
\caption{Pairwise Simpson tensions across the three production
rows.}
\label{tab:pairwise}
\end{table}

All three tensions lie below $0.2\sigma$. The combined uncertainties
are computed directly from the row-level standard errors as
$\sigma_{\rm pair} = \sqrt{\sigma_i^2 + \sigma_j^2}$.

\subsection{Interpretation: stability across the measured range}
\label{subsec:lattice-independence}

The three production rows differ simultaneously in lattice volume
and in the caloron scale parameter $\rho$ quoted in lattice units.
They therefore probe distinct combinations of finite-volume and
discretization effects. The observed stability of
$\delta_\xi \log\zeta$ at sub-$0.2\sigma$ across this range is the
direct empirical content of the present section. The statement is stability across the measured range, not a
formal continuum extrapolation; we do not infer $a \to 0$ or
$V \to \infty$ from three data points.

The tightness of the pairwise tensions at the noise level of
Table~\ref{tab:production-results} reflects the paired-difference
construction of Section~\ref{subsec:simpson-setup}: the low-variance
component of each Simpson estimator is the paired difference
$h(\xi_j) - h(\xi_0)$ computed against a common noise seed within
each row, so that the Monte Carlo variance of the Simpson integrand
is dominated by the difference rather than the absolute node values.
The resulting row-level standard errors $\sigma_i$ of
Table~\ref{tab:production-results} translate directly into the tight
pairwise tensions of Table~\ref{tab:pairwise}, and the
inverse-variance weighted three-row result of~\eqref{eq:bracket-result}
enters the assembled result of Section~\ref{sec:assembled}.


\section{Assembled result and error budget}
\label{sec:assembled}

\subsection{Summary of factors}
\label{subsec:assembly-summary}

Table~\ref{tab:assembly} collects the quantities that enter the
assembled result directly. Each is quoted at the precision
established in the section where it is computed; no quantity is
re-measured here.

\begin{table}[h]
\centering
\begin{tabular}{lll}
\toprule
Quantity & Value & Reference \\
\midrule
$e^{-S_{\rm inst}}$
  & $\sim 10^{-120}$ \,(conditional on $S_{\rm inst} \simeq 276$)
  & \cite{companion-structural} \\
$\Jphys$
  & $+0.118164$ \,(with $0.4\%$ row-to-row variation)
  & Section~\ref{sec:jacobian} \\
$\zeta(0)$
  & $8.42 \pm 0.08$ \,(Row A; $0.4\%$ row-to-row)
  & Eq.~\eqref{eq:zeta0-value} \\
$\Bobs \equiv 1 - \zeta(X_*)/\zeta(0)$
  & $0.811 \pm 0.023$ \,(three-row weighted)
  & Eq.~\eqref{eq:bracket-result} \\
$\cos\Theta_\gamma$
  & $+1$ \,(parity-even class)
  & Theorem~\ref{thm:phase} \\
\bottomrule
\end{tabular}
\caption{Quantities entering the assembled result. The ordering
matches the factor sequence of the master
formula~\eqref{eq:master}.}
\label{tab:assembly}
\end{table}

The adaptive-prime prescription and the projected-complement
exact-prime cross-check are validation items for the branch-response
determinant; they are not additional multiplicative factors in the
master formula.

\subsection{Assembly arithmetic}
\label{subsec:assembly-arithmetic}

For the numerical assembly it is convenient to collect the
lattice-determined prefactor multiplying $e^{-S_{\rm inst}}$ into a
single quantity,
\begin{equation}
\mathcal{K}
\;\equiv\;
\Jphys\,\bigl[\zeta(0) - \zeta(X_*)\bigr]\,\cos\Theta_\gamma .
\label{eq:K-def}
\end{equation}
In the normalization convention used throughout the production
determinant, the measured Jacobian, phase, and branch-response
input combine to
\begin{equation}
\mathcal{K} \;=\; \Jphys \cdot \zeta(0) \cdot \Bobs \cdot \cos\Theta_\gamma
\;=\; (0.118164)(8.42 \pm 0.08)(0.811 \pm 0.023)(+1)
\;=\; 0.807 \pm 0.024.
\label{eq:K-value}
\end{equation}
Conditional on $S_{\rm inst} \simeq 276$, the assembled result is
therefore
\begin{equation}
\boxed{\;
\frac{\Lobs}{\MP^4}
\;=\;
(0.81 \pm 0.03) \times 10^{-120}
\quad
\text{conditional on}\ S_{\rm inst} \simeq 276.
\;}
\label{eq:assembled-result}
\end{equation}

\subsection{Conditional scope}
\label{subsec:assembly-conditional}

The assembled result~\eqref{eq:assembled-result} is conditional on
two inputs external to the direct lattice measurement reported
here. First, the hidden-sector action is taken as
$S_{\rm inst} \simeq 276$ from a separate
analysis~\cite{companion-structural}; all numerical content of this
paper is independent of the hidden-sector algebra, $S_{\rm inst}$
enters the assembled result only multiplicatively, and its
sensitivity is quantified in Appendix~\ref{app:S-sensitivity}.
Second, the branch-response determinant is computed using the
adaptive-prime prescription of Section~\ref{sec:spectral},
validated by the three independent checks
of~Section~\ref{subsec:operational-closure}. A first-principles dressed
analytic basis for the adaptive prime is not required for the
present result and is the subject of separate technical work.

\subsection{Error budget}
\label{subsec:error-budget}

The uncertainty on the lattice prefactor $\mathcal{K}$
of~\eqref{eq:K-value} is dominated by the three-row weighted
bracket uncertainty of Appendix~\ref{app:error-budget}.
Table~\ref{tab:error-budget} collects the contributions at the
level resolved in the present paper.

\begin{table}[h]
\centering
\begin{tabular}{lll}
\toprule
Source & Relative size & Contribution to $\sigma(\mathcal{K})$ \\
\midrule
Bracket uncertainty
  & $0.023 / 0.811 \simeq 2.8\%$
  & $2.3 \times 10^{-2}$ \\
$\zeta(0)$ normalization
  & $0.08 / 8.42 \simeq 1.0\%$
  & $8 \times 10^{-3}$ \\
Jacobian row-to-row variation
  & $0.4\%$
  & $3 \times 10^{-3}$ \\
Phase
  & exact in the parity-even class
  & $0$ \\
Sub-dominant numerical systematics
  & Appendix~\ref{app:error-budget}
  & $< 10^{-3}$ \\
\midrule
Quadrature sum
  & bracket-dominated
  & $2.5 \times 10^{-2}$ \\
\bottomrule
\end{tabular}
\caption{Contributions to the uncertainty on the prefactor
$\mathcal{K}$. The boxed result~\eqref{eq:assembled-result} quotes
$\pm 0.03$, a conservative rounding of the quadrature sum $0.025$
reserving margin for higher-order systematics not quantified at the
current precision, including residual
contributions from the anchor-branch normalization entering at the
$1\%$ level.}
\label{tab:error-budget}
\end{table}

The quoted $\pm 0.03$ in~\eqref{eq:assembled-result} is therefore a
bracket-dominated uncertainty, with all other identified numerical
effects subdominant at the present precision.

\subsection{Named results of this paper}
\label{subsec:named-results}

The logical spine of the computation consists of three named
results:
\begin{itemize}
\item Proposition~\ref{prop:factor-partition}
(Section~\ref{subsec:framework-master}): the master formula factorizes
into UV-finite quantities after the analytic constitutive
subtraction.
\item Theorem~\ref{thm:phase} (Section~\ref{sec:phase}): in the
parity-even class, $\cos\Theta_\gamma = +1$ on $\bR^4$,
$S^1 \times \bR^3$, and $T^4$.
\item Proposition~\ref{prop:universality}
(Section~\ref{subsec:universality}): across three production lattices
spanning $11\times$ in volume, the Simpson nonlinear correction
remains in the range $\Delta_{\rm NL} \in [13.5\%,\,16.4\%]$ with
unweighted mean $14.6\%$.
\end{itemize}
In addition, Section~\ref{subsec:cross-check-exact} reports a direct
cross-prescription measurement: on Row~A with $N = 800$ paired
samples, the adaptive prime and the projected-complement exact
prime agree at $0.77\sigma$. This is a prescription-level
consistency check, not a separate named theorem.

\section{Discussion}
\label{sec:discussion}

\subsection{Conditional scope}
\label{subsec:discussion-scope}

The assembled value~\eqref{eq:assembled-result} is a numerical
result conditional on two structural inputs. First, the
hidden-sector action $S_{\rm inst} \simeq 276$ is taken as input
from the companion analysis~\cite{companion-structural} and is not
recomputed in the present paper; the sensitivity of the assembled
value to this input is quantified in
Appendix~\ref{app:S-sensitivity}. Second, the branch-response
determinant is computed under the adaptive-prime prescription of
Section~\ref{sec:spectral}, whose validity at the level of the
final bracket is supported by the three independent checks
of~Section~\ref{subsec:operational-closure}. Neither conditional is an
isolated empirical assumption: each is framework-level structural
input derived in~\cite{companion-structural}, and the paper's
numerical content is conditional on that framework as a whole. The
claim of the paper is therefore structural and numerical rather
than empirical in isolation: within the stated framework, the
lattice computation yields a cosmological constant of the observed
order with prefactor $0.81 \pm 0.03$.

\subsection{Gravitational-only coupling of the hidden sector}
\label{subsec:discussion-coupling}

The master formula~\eqref{eq:master} is derived under the
assumption that the hidden sector generating the instanton action
couples to the visible sector only through the emergent metric
$g_{\rm eff}$. Direct couplings (through scalar labels, gauge
portals, or kinetic mixing) would introduce additional operator
content that is not included in the lattice pipeline of the present
paper. Under gravitational-only coupling, the only visible-sector
consequence of the hidden sector is the cosmological-constant value
itself: the hidden sector does not produce an independent
particle-physics or cosmological signature at the energies probed
by current experiment. The structural derivation of the
gravitational-only requirement is given
in~\cite{companion-structural}.

\subsection{Observational posture}
\label{subsec:discussion-observation}

The empirical anchor of the present computation is the magnitude
of the cosmological constant it reproduces. Under the
gravitational-only coupling of Section~\ref{subsec:discussion-coupling},
we do not predict an independent empirical signature beyond the
cosmological-constant magnitude, and we do not discuss here the
downstream cosmological consequences (such as contributions to
$\Delta N_{\rm eff}$, gravitational-wave polarization content, or
gravitational-slip corrections) that would follow from a specific
completion of the hidden sector. A systematic analysis of such
downstream signatures requires a separate bridge between the
framework~\cite{companion-structural} and the precision-cosmology
data and is deferred to subsequent work.

\subsection{Relation to existing approaches}
\label{subsec:discussion-comparison}

The construction differs in character from several established
approaches to the cosmological constant problem.
\begin{itemize}
\item \emph{Supersymmetric cancellation}~\cite{Martin12}. The
present paper does not rely on supersymmetry. Polynomial
ultraviolet hierarchies are removed by the analytic constitutive
subtraction of~\cite{companion-structural}, not by
boson--fermion cancellation; a residual nonperturbative
contribution is then computed explicitly.
\item \emph{Anthropic or landscape selection}~\cite{Weinberg89}.
The value~\eqref{eq:assembled-result} is derived within a specific
framework rather than selected from a distribution of vacua.
\item \emph{Scalar-only local transmutation.}
The present paper does not invoke a local scalar-adjustment
mechanism. Within the framework used here, scalar-only
transmutation of the vacuum energy at the observed magnitude is
excluded by the no-go results
of~\cite{companion-structural}.
\item \emph{Volovik-type superfluid-vacuum
analogies}~\cite{Volovik03}. The present construction sits squarely within the induced-gravity
tradition that includes these analogies, and differs from them in
that the induced effective action is supplied concretely by the
proto-field sector~\cite{companion-structural} and the
nonperturbative residual is computed on the lattice rather than
extracted from an analogue.
\end{itemize}

\subsection{Open technical items}
\label{subsec:discussion-open}

Three items remain open at the technical rather than architectural
level. First, a first-principles dressed analytic basis for the
adaptive-prime prescription would replace the IPR-based
identification of $V_{20}$ by a closed-form construction capturing
the compact-$T^4$ dressing of the moduli-space tangents; the
adaptive-prime prescription used in the present paper is validated
by the three independent checks
of~Section~\ref{subsec:operational-closure} without relying on such a
basis. Second, an analytic upgrade of the subdominant determinant
coefficient entering the hidden-sector computation
of~\cite{companion-structural} would replace the present numerical
value by a closed-form determination; this does not affect the
lattice result reported here. Third, a first-principles antecedent
for the hidden-sector action value $S_{\rm inst} \simeq 276$ is
the subject of separate work within the same
framework~\cite{companion-structural}; the numerical sensitivity
of the assembled result to this input is quantified in
Appendix~\ref{app:S-sensitivity}. None of these items is required
for the numerical content of the present paper.

\section{Conclusion}
\label{sec:conclusion}

Within the coframe--teleparallel induced-gravity framework
of~\cite{companion-structural}, the nonperturbative residual of
the cosmological constant problem has been computed on the
lattice. The four UV-finite factors of the master formula have
been evaluated independently, and the assembled result is
\[
\boxed{\;
\frac{\Lobs}{\MP^4}
\;=\;
(0.81 \pm 0.03) \times 10^{-120}
\qquad\text{conditional on}\ S_{\rm inst} \simeq 276.
\;}
\]
Beyond this number, the computation establishes two structural
findings. The Simpson nonlinear correction $\Delta_{\rm NL}$ is a
stable prediction of the branch-bias mechanism, measured in the
range $[13.5\%,\,16.4\%]$ across three lattices spanning $11\times$
in volume; and the parity-even phase $\cos\Theta_\gamma = +1$ is a
theorem on every background geometry considered. The production
prescription is additionally supported by the cross-prescription
agreement of~Section~\ref{subsec:cross-check-exact}, which finds
$0.77\sigma$ tension between the adaptive prime and the
projected-complement exact prime on Row~A. The remaining open
technical items --- a dressed first-principles basis for the
adaptive-prime prescription, an analytic upgrade of a subdominant
determinant coefficient entering the hidden-sector computation,
and a first-principles antecedent for the value of
$S_{\rm inst}$ --- are technical, not architectural, and are
addressed in separate companion work within the same framework.



\appendix

\section{Numerical conventions and code availability}
\label{app:conventions}

This appendix collects the conventions and numerical parameters held fixed
throughout the lattice computation. Its purpose is to render the production
pipeline reproducible from the equations of the main text. Wherever a
convention has an obvious alternative (for example, an overall operator
normalisation), we state the choice explicitly; the numerical values
reported in the main text correspond exclusively to the conventions listed
here.

\subsection{Lattice geometry and index layout}
\label{app:conv-geometry}

We work on a periodic four-dimensional hypercubic lattice of size
$L_t \times L_s^3$, with $L_t$ the temporal extent and $L_s$ the common
spatial extent. The lattice spacing $a$ is used as the unit of length: all
dimensional quantities are expressed in lattice units unless stated
otherwise. The temporal direction is $\mu = 0$ and the spatial directions
are $\mu = 1, 2, 3$; periodic boundary conditions are imposed in every
direction.

Gauge links $U_{x,\mu} \in SU(2)$ are stored in a rank-seven complex array
of shape $(L_t,\,L_s,\,L_s,\,L_s,\,4,\,2,\,2)$, corresponding to the
coordinate $(x_0,x_1,x_2,x_3,\mu)$ and the two colour indices. Unitarity of
each link is enforced at assembly and monitored throughout the run; the
absolute deviation $\|U^\dagger U - \mathbb{1}\|_{\max}$ stays below
$10^{-15}$ on every production background.

A lattice adjoint-vector field $v^a_\mu(x)$ is stored as a flat complex
array $v \in \mathbb{C}^{12 V}$, $V \equiv L_t L_s^3$, under the linear
index layout
\begin{equation}
v[\,12\,x + 3\mu + a\,] \;=\; v^a_\mu(x),
\qquad
x = 0,\dots,V-1,\;\;\mu = 0,\dots,3,\;\;a = 1,2,3.
\label{eq:app-index-layout}
\end{equation}
All sparse matrices representing lattice operators on this space have
shape $(12 V) \times (12 V)$ and are ordered according
to~\eqref{eq:app-index-layout}.

The Lie algebra $\mathfrak{su}(2)$ is represented by the anti-Hermitian
generators $T^a = \tfrac{i}{2}\sigma^a$, so that $T^{a\dagger} = -T^a$, with
$\sigma^a$ the Pauli matrices; we use the trace normalisation
$\Tr(T^a T^b) = -\tfrac{1}{2}\delta^{ab}$. Adjoint indices $a = 1, 2, 3$
are raised and lowered with $\delta^{ab}$.

\subsection{Anisotropic metric and the branch parameter}
\label{app:conv-metric}

The branch deformation of the main text is implemented at the level of the
background metric by anisotropic stretching of the temporal direction,
\begin{equation}
g_{\mu\nu}(\xi) \;=\; \mathrm{diag}(\xi^2,\,1,\,1,\,1),
\qquad
\xi \in [1,\,1+\epsilon],
\label{eq:app-aniso-metric}
\end{equation}
with $\epsilon = 0.09$ the disformal parameter held fixed throughout. The
reference value $\xi = 1$ corresponds to the undeformed metric underlying
the Harrington--Shepard caloron background. The three-point Simpson
quadrature used in Section~\ref{sec:finite-xi} evaluates the branch
response at
\begin{equation}
\xi \;\in\; \{\,1.000,\;1.045,\;1.090\,\},
\label{eq:app-xi-grid}
\end{equation}
corresponding to $\xi = 1$, $\xi = 1 + \epsilon/2$, and $\xi = 1 + \epsilon$.

\subsection{Lattice operators}
\label{app:conv-operators}

Four lattice operators enter the branch-response computation. We record
here only the sign conventions and regulator choices actually used in the
production code. A common infrared regulator $m^2 = 0.01$ is added to the
Laplace-type blocks throughout.

\paragraph{Adjoint vector operator.}
The vector operator used for the first-order response and the Simpson
integrand is
\begin{equation}
\Delta_1 \;=\; -D^2 \,+\, 2F \,+\, m^2\,\mathbb{1},
\label{eq:app-Delta1}
\end{equation}
where $D^2$ is the adjoint gauge-covariant Laplacian and $F$ denotes the
adjoint field-strength insertion,
\begin{equation}
(F v)^a_\mu(x)
\;=\;
2\,\bigl[F_{\mu\nu}(x)\bigr]^a_{\phantom{a}b}\,v^b_\nu(x),
\label{eq:app-F-insertion}
\end{equation}
with $F_{\mu\nu}$ the clover-average field-strength tensor on the adjoint
bundle, and the sum over $\nu$ understood. The overall factor of $2$ is
the convention implemented in
\texttt{src/spectral.py:build\_vector\_operator}.

\paragraph{Adjoint scalar Laplacian.}
The scalar Laplacian,
\begin{equation}
\Delta_0 \;=\; -D^2 \,+\, m^2\,\mathbb{1},
\label{eq:app-Delta0}
\end{equation}
appears in the longitudinal-mode subtraction of the branch response and is
implemented as \texttt{src/spectral.py:build\_adjoint\_laplacian}.

\paragraph{Temporal operator.}
The temporal-only adjoint Laplacian,
\begin{equation}
K_0 \;=\; -D_0^2,
\label{eq:app-K0}
\end{equation}
is the building block of the anisotropic term $A'$ defined below and is
obtained from the same routine with the spatial covariant derivatives
disabled.

\paragraph{Anisotropic insertion and $\delta F$.}
The derivative of $\Delta_1$ along the branch parameter combines the
temporal block of the Laplacian and the $\xi$-variation of the field
strength,
\begin{equation}
A' \;\equiv\; \partial_\xi \Delta_1(\xi)\bigr|_{\xi=1}
\;=\; +\,2\,K_0 \,+\, 2\,\delta F,
\label{eq:app-Aprime}
\end{equation}
where the $\delta F$ insertion is built by
\texttt{build\_delta\_F\_operator(U, dims)} with
\texttt{dims = (L\_t, L\_s, L\_s, L\_s)}, and the sign in
\eqref{eq:app-Aprime} is the one consistent with the metric
stretch~\eqref{eq:app-aniso-metric}. All numerical values reported in this
paper use the sign convention of~\eqref{eq:app-Aprime}. The
combination~\eqref{eq:app-Aprime} is the anisotropic datum evaluated at
each Simpson node.

\paragraph{Gauge-cleaning projector.}
The raw left-trivialized tangents of the collective coordinates, computed
by finite differences in Section~\ref{sec:jacobian}, contain residual
gauge components. These are removed by the projector
$P_\perp = \mathbb{1} - V_{\rm gauge} V_{\rm gauge}^{\!\top}$, where
$V_{\rm gauge}$ is an orthonormal basis of the gauge-orbit tangents at the
caloron configuration. This cleaning is performed by
\texttt{jacobian\_exact.py:gauge\_clean\_tangent}; the residual gauge
overlap after cleaning, measured through the orientation indicator
$G_{\rm orient}$, is of order $4\times 10^{-6}$ on every production
background.

\subsection{Free-background subtraction and the Jacobian
contribution to $\delta_\xi \log\zeta$}
\label{app:conv-subtraction}

The branch-response quantity $\delta_\xi \log\zeta$ reported in
the main text is the fully subtracted quantity
\begin{equation}
\delta_\xi \log\zeta
\;=\;
\bigl[\log\zeta_{\rm inst}(\xi) - \log\zeta_{\rm inst}(1)\bigr]
\;-\;
\bigl[\log\zeta_{\rm free}(\xi) - \log\zeta_{\rm free}(1)\bigr],
\label{eq:app-delta-subtracted}
\end{equation}
where $\log\zeta$ receives contributions from the operator
determinants, the ghost sector, and the collective-coordinate
Gram. The Gram contribution on the instanton background is
$\tfrac{1}{2}\log\det G_5 = \log\Jphys$; on the free background
the Gram does not appear (there are no collective coordinates).
Under~\eqref{eq:app-delta-subtracted}, the Gram contributes
$\log\Jphys$ at $\xi = 1$ and, because the cleaned left-trivialized
tangents of Section~\ref{sec:jacobian} are $\xi$-independent under
the metric stretch~\eqref{eq:app-aniso-metric}, the Gram
contribution is likewise $\xi$-independent. It therefore enters
$\delta_\xi \log\zeta$ as an additive constant, equal numerically
to $\Jphys$ of~\eqref{eq:Jphys-value} under the FD-step
normalisation of Table~\ref{tab:app-parameters}. This fixes the
identification $J_{\rm anchor} = \Jphys$ used in the Simpson
formula~\eqref{eq:simpson-quadrature}. The equality is
convention-dependent through the choice of FD step and the
placement of normalisation factors; the assembled
$\delta_\xi \log\zeta$ is invariant under these choices.

\subsection{Stochastic tracing and solvers}
\label{app:conv-stochastic}

Traces of the form $\Tr\bigl(P\,\mathcal{O}(\xi)\,P\bigr)$ with $P$ a
projector are estimated by paired Hutchinson noise using $N_S$ Rademacher
vectors at a common random seed across all values of
$\xi$~\cite{Hutchinson1990}. 
The paired construction ensures that the difference
$\Tr_P(\xi_1) - \Tr_P(\xi_2)$ is computed against a common noise
realisation, so that the Monte Carlo variance of the Simpson integrand is
dominated by the difference rather than by the absolute values. In
production we use
\begin{equation}
N_S \;=\; 80,\qquad \text{seed}\;=\;42,
\label{eq:app-seed}
\end{equation}
and the paired differences $h(1.045) - h(1.000)$ and $h(1.090) - h(1.000)$
have standard errors an order of magnitude below those of the absolute
integrand values $h(\xi)$ themselves.

Linear solves against $\Delta_1$ on the projected complement of the
adaptive-prime subspace are performed by conjugate gradient. The
production solver is the GPU-resident routine
\texttt{cupyx.scipy.sparse.linalg.cg} wrapping a CuPy
\texttt{LinearOperator} that implements
\eqref{eq:app-Delta1}--\eqref{eq:app-F-insertion} directly in device
memory. The solver tolerance is $10^{-10}$ in relative residual, and
convergence was uniform on every production row.

\subsection{Physical parameters and reference values}
\label{app:conv-parameters}

Table~\ref{tab:app-parameters} collects the parameters held fixed across
every production computation in the main text.

\begin{table}[h]
\centering
\begin{tabular}{lll}
\toprule
Quantity & Value & Role \\
\midrule
$m^2$             & $0.01$                          & infrared regulator, all Laplacians \\
$\epsilon$        & $0.09$                          & disformal branch parameter \\
$\xi$ grid        & $\{1.000,\,1.045,\,1.090\}$     & Simpson quadrature nodes \\
$N_S$             & $80$                            & paired Hutchinson samples \\
seed              & $42$                            & common noise seed across $\xi$ \\
CG tolerance      & $10^{-10}$                      & relative residual, projected solve \\
FD step (transl.) & $\delta = 0.1$                  & translation tangents \\
FD step (scale)   & $\delta = 0.05$                 & scale tangent \\
\bottomrule
\end{tabular}
\caption{Numerical parameters held fixed across every production row in
the main text. Finite-difference (FD) steps are those used for the
left-trivialized tangents of Section~\ref{sec:jacobian}.}
\label{tab:app-parameters}
\end{table}

\subsection{Software stack and code availability}
\label{app:conv-software}

The production pipeline is written in Python~3.11, with NumPy and SciPy
providing the CPU linear algebra~\cite{Harris2020,Virtanen2020}
and CuPy providing the GPU linear algebra~\cite{Okuta2017}.
Sparse matrices are stored in CSR format on the CPU and in the
corresponding CuPy CSR format on the GPU. The only code modification from
the CPU development pipeline to the H100 production pipeline is the
replacement of the SciPy projected-CG solver by the CuPy projected-CG
solver; all operator definitions and all tangent computations are shared.

All production scripts and the supporting \texttt{src/} module
library are collected in a single repository. The core production
scripts are: \texttt{gate2\_runpod\_h100\_v2.py} (transported-prime
Simpson campaign on the three production rows);
\texttt{ipr\_gated\_prime.py} (adaptive-prime spectral selection
and $\delta_\xi \log\zeta$ evaluation at fixed~$\xi$);
\texttt{exact\_prime\_projected.py} (projected conjugate-gradient
solver); \texttt{exact\_prime\_F\_G.py} (the
\texttt{build\_delta\_F\_operator} routine);
\texttt{jacobian\_left\_triv.py} (five left-trivialized tangents
and the $5 \times 5$ Gram matrix);
\texttt{jacobian\_exact.py} (gauge cleaning);
\texttt{verify\_jacobian\_infra.py} (infrastructure self-tests
of Appendix~\ref{app:infra-tests}); and
\texttt{mode\_diagnostic.py} (spectral classification used in
Section~\ref{sec:spectral}).



\section{Transported-prime algorithm and validation}
\label{app:transported-prime}

The finite-$\xi$ Simpson estimator of Section~\ref{sec:finite-xi} requires
the adaptive-prime subspace $V_{20}$ to be defined consistently at each of
the three quadrature nodes $\xi \in \{1.000,\,1.045,\,1.090\}$.
Reconstructing $V_{20}$ from scratch at every node can introduce a
discontinuity in the projector when mode identities cross between nodes;
such a discontinuity contaminates the paired-difference structure on which
the low-variance property of the Simpson estimator depends. The
transported-prime prescription adopted in production avoids this
contamination by anchoring the prime at $\xi_0 = 1$ and re-projecting onto
the $\Delta_1(\xi_j)$ eigenbasis at every other node via greedy overlap
matching. This appendix specifies the algorithm, records its role
relative to first-order parallel transport, and reports the overlap
diagnostics on Row~A.

\subsection{Algorithm}
\label{app:transport-algorithm}

Let $V_{20}(\xi_0) = \mathrm{span}\{\psi_1(\xi_0),\dots,\psi_{20}(\xi_0)\}$
denote the anchor subspace at $\xi_0 = 1$, identified once by direct IPR
gating on the $\Delta_1(\xi_0)$ eigenvectors
(Section~\ref{sec:spectral}). At each subsequent Simpson node
$\xi_j \in \{1.045,\,1.090\}$, the transported prime $V_{20}(\xi_j)$ is
constructed by the following procedure.

\begin{quote}
\textbf{Algorithm B.1}~(\emph{transported prime at $\xi_j$}).\medskip

\noindent\textit{Inputs.} Anchor basis
$\{\psi_\alpha(\xi_0)\}_{\alpha=1}^{20}$; transported operator
$\Delta_1(\xi_j)$; candidate pool size $N_c > 20$.

\noindent\textit{Output.} Transported basis
$\{\hat\psi_\alpha(\xi_j)\}_{\alpha=1}^{20}$ spanning $V_{20}(\xi_j)$.

\begin{enumerate}
\item Compute the $N_c$ smallest-eigenvalue eigenpairs
$\{(\lambda_k(\xi_j),\,\phi_k(\xi_j))\}_{k=1}^{N_c}$ of $\Delta_1(\xi_j)$
by shift-invert Lanczos, returning eigenvectors orthonormal to machine
precision.

\item Form the overlap matrix
\[
O_{\alpha k}
\;=\;
\bigl|\langle \psi_\alpha(\xi_0) \,|\, \phi_k(\xi_j) \rangle\bigr|,
\qquad
\alpha = 1,\dots,20,\;\; k = 1,\dots,N_c.
\]

\item Initialise assigned sets $\mathcal{A} = \emptyset$ and candidate
pool $\mathcal{C} = \{1,\dots,N_c\}$.

\item \textbf{For} $n = 1,\dots,20$:
\begin{enumerate}
\item[(a)] find the maximising pair
$(\alpha^*,\,k^*) \,=\, \arg\max_{\alpha\notin\mathcal{A},\;k\in\mathcal{C}}
O_{\alpha k}$;
\item[(b)] set $\hat\psi_{\alpha^*}(\xi_j) \leftarrow \phi_{k^*}(\xi_j)$;
\item[(c)] update $\mathcal{A} \leftarrow \mathcal{A} \cup \{\alpha^*\}$,
$\mathcal{C} \leftarrow \mathcal{C}\setminus\{k^*\}$.
\end{enumerate}

\item Return $\{\hat\psi_\alpha(\xi_j)\}_{\alpha=1}^{20}$.
\end{enumerate}
\end{quote}

The output is orthonormal by construction, because the $\phi_k(\xi_j)$ are
exact eigenvectors of the Hermitian operator $\Delta_1(\xi_j)$; no
Gram--Schmidt step is required. In production, $N_c$ is chosen as the
smallest Lanczos window for which the twenty transported matches are all
contained in the returned spectrum and the greedy assignment is unique.
This criterion, rather than a fixed numerical value of $N_c$, is the
reproducibility condition used throughout. The greedy rule is
one-to-one: each anchor mode is paired with a distinct transported
eigenvector.

Equivalently, the transported prime used in the trace computation is the
rank-$20$ projector
\begin{equation}
P_{20}(\xi_j)
\;=\;
\sum_{\alpha=1}^{20}
|\hat\psi_\alpha(\xi_j)\rangle\!\langle\hat\psi_\alpha(\xi_j)|,
\label{eq:app-P20}
\end{equation}
and the production trace is evaluated on its orthogonal complement.

\subsection{Re-projection versus first-order parallel transport}
\label{app:transport-vs-fopt}

Algorithm~B.1 may be compared with first-order parallel transport (FOPT),
in which each anchor mode is evolved by the Rayleigh--Schr\"odinger
linearisation at $\xi_0$. For the present paper the relevant comparison
is not between two approximate transports, but between an exact
re-projection onto the eigenspace of $\Delta_1(\xi_j)$ and a first-order
truncation of that eigenspace.

\begin{lemma}[Why exact re-projection is used]
\label{lem:rep-vs-fopt}
Assume that, at a node $\xi_j$, the twenty-mode cluster selected by the
IPR criterion remains isolated from the bulk and that the greedy overlap
assignment of Algorithm~B.1 is unique. Then:
\begin{enumerate}
\item[(i)] the basis returned by Algorithm~B.1 consists of exact
eigenvectors of $\Delta_1(\xi_j)$;
\item[(ii)] its span equals the transported low-manifold $V_{20}(\xi_j)$,
up to basis rotations inside any exactly degenerate eigenspace;
\item[(iii)] a first-order parallel-transport basis reproduces the same
eigenspace only up to the truncation error of its linear expansion in
$\xi_j - \xi_0$.
\end{enumerate}
\end{lemma}

\begin{proof}
Statement (i) is immediate from step~1 of Algorithm~B.1: each selected
$\hat\psi_\alpha(\xi_j)$ is one of the Lanczos eigenvectors
$\phi_k(\xi_j)$ and therefore satisfies
\[
\Delta_1(\xi_j)\,\hat\psi_\alpha(\xi_j)
\;=\;
\lambda_k(\xi_j)\,\hat\psi_\alpha(\xi_j).
\]
For (ii), the isolation of the twenty-mode cluster and uniqueness of the
greedy overlap assignment imply that the twenty selected eigenvectors
are precisely the twenty eigenvectors defining the low-manifold at
$\xi_j$. Any residual ambiguity inside an exactly degenerate eigenspace
is basis-dependent but projector-independent. For (iii), the FOPT
construction truncates the $\xi$-dependence of the eigenvectors at first
order,
\begin{equation}
\hat\psi_\alpha^{\rm FOPT}(\xi_j)
\;=\;
\psi_\alpha(\xi_0)
\;+\;
(\xi_j - \xi_0)\!\!\sum_{\beta \neq \alpha}
\frac{\langle\psi_\beta(\xi_0)\,|\,A'\,|\,\psi_\alpha(\xi_0)\rangle}
     {\lambda_\alpha(\xi_0) - \lambda_\beta(\xi_0)}\,
\psi_\beta(\xi_0),
\label{eq:app-FOPT}
\end{equation}
with $A' = \partial_\xi \Delta_1(\xi)\bigr|_{\xi=\xi_0}$
of~\eqref{eq:app-Aprime}; so its span differs from the exact eigenspace
by the omitted higher-order terms in the same expansion. Exact
re-projection removes that transport-truncation error by reconstructing
the eigenspace directly at each Simpson node.
\end{proof}

This is the only sense in which Algorithm~B.1 is stronger than FOPT in
the present paper: it replaces a first-order approximation to the
deformed eigenspace by an exact eigenspace reconstruction at each
quadrature node.

\subsection{Overlap diagnostics on Row A}
\label{app:transport-overlaps}

The quality of the transport is monitored by the mode-wise minimum
overlap
\begin{equation}
\Omega_{\min}(\xi_j)
\;\equiv\;
\min_{\alpha = 1,\dots,20}\,
\bigl|\langle\psi_\alpha(\xi_0)\,|\,\hat\psi_\alpha(\xi_j)\rangle\bigr|,
\label{eq:app-Omega-min}
\end{equation}
which equals unity if anchor and transported bases agree mode by mode
and degrades toward zero in the presence of mode crossings or an
ambiguous greedy assignment. Figure~\ref{fig:app-overlaps} records the
mode-wise overlaps on Row~A, which serves as the representative
transport diagnostic for the production campaign.

\begin{figure}[h]
\centering
\includegraphics[width=0.65\textwidth]{fig_transport_overlaps_rowA.pdf}
\caption{Transport diagnostics for Algorithm~B.1 on Row~A.
Panel~(a): per-mode assigned overlaps
$|\langle\psi_\alpha(\xi_0)\,|\,\hat\psi_\alpha(\xi_j)\rangle|$
at the two non-anchor Simpson nodes $\xi_1 = 1.045$ and
$\xi_2 = 1.090$, with $\Omega_{\min} \in [0.57,\,0.61]$ reflecting
degenerate-subspace rotation inside Tier~1 and Tier~2 (permitted by
Lemma~\ref{lem:rep-vs-fopt}(ii)). Panel~(b): the subspace overlap
$\mathcal{S}(\xi) \equiv \mathrm{Tr}[P_{V_{20}}(\xi_0)\,P_{V_{20}}(\xi)]
/ N_{\rm anchor}$, invariant under unitary rotations within degenerate
eigenspaces. The measured values $\mathcal{S}(1.045) = 0.99$ and
$\mathcal{S}(1.090) = 0.96$ confirm that the twenty-dimensional low
manifold transports as a coherent subspace across the Simpson window,
validating the premise of Algorithm~B.1.}
\label{fig:app-overlaps}
\end{figure}

On Row~A the greedy assignment is unique in every case, and the
subspace-overlap diagnostic $\mathcal{S}(\xi)$ (Fig.~\ref{fig:app-overlaps}
panel~b) exceeds $0.96$ at both non-anchor nodes, confirming that
$V_{20}(\xi_0)$ and $V_{20}(\xi_j)$ span essentially the same
twenty-dimensional subspace of $\mathbb{C}^{12V}$. The per-mode
assigned overlaps (panel~a) dip to $\Omega_{\min} \simeq 0.57$ at
$\xi = 1.045$ and $\Omega_{\min} \simeq 0.61$ at $\xi = 1.090$,
reflecting unitary rotations inside the exactly- and
near-degenerate eigenspaces of Tier~1 and Tier~2 (Lemma~B.1~(ii)
explicitly permits this). This is sufficient for the
transported-prime construction used in
Section~\ref{sec:finite-xi}: what the Simpson estimator requires is
subspace preservation of $V_{20}$, not preservation of individual
mode identities within degenerate subspaces.

\subsection{Computational cost}
\label{app:transport-cost}

The transported-prime construction is subdominant to the projected-CG
trace estimation in wall time. The Lanczos solve at each of the two
non-anchor nodes computes $N_c$ eigenpairs of $\Delta_1(\xi_j)$ to a
relative residual of $10^{-8}$ on each eigenpair, and the greedy matching
itself is an $\mathcal{O}(N_c \cdot 20)$ operation on a small overlap
matrix. On the H100 pipeline, the combined cost of Algorithm~B.1 at both
non-anchor nodes amounts to a few minutes per production row and is
subdominant to the projected-CG trace, which dominates the wall time of
the Simpson campaign.



\section{Infrastructure verification tests}
\label{app:infra-tests}

The lattice pipeline includes a compact self-test suite,
\texttt{verify\_jacobian\_infra.py}, whose purpose is to guard against
silent regressions in the lattice primitives, the operator constructors,
the gauge-cleaning projector, and the $5\times 5$ Gram assembly. The
suite exercises the code paths on which the numerical content of
Section~\ref{sec:jacobian} and Section~\ref{sec:finite-xi} depends; it
is a software-integrity check rather than a physics validation.
Table~\ref{tab:app-infra} summarises the version used for the production
runs reported in the main text.

\begin{table}[h]
\centering
\begin{tabular}{clll}
\toprule
\# & Test & Primitive exercised & Pass criterion \\
\midrule
1 & Index round-trip
  & \texttt{v[12x + 3$\mu$ + a]} layout
  & bijection on $12V$ entries \\
2 & Link unitarity
  & SU(2) link assembly
  & $\|U^\dagger U - \mathbb{1}\|_{\max} < 10^{-15}$ \\
3 & Clover field strength
  & $F_{\mu\nu}$ adjoint projection
  & anti-Hermiticity $\leq 10^{-14}$ \\
4 & Vector operator sign check
  & $\Delta_1 = -D^2 + 2F + m^2$
  & lowest computed Ritz value $\geq m^2$ \\
5 & Scalar Laplacian sign check
  & $\Delta_0 = -D^2 + m^2$
  & lowest computed Ritz value $\geq m^2$ \\
6 & Temporal Laplacian
  & $K_0 = -D_0^2$ block consistency
  & Hermiticity $\leq 10^{-14}$ \\
7 & Free-field tangent recovery
  & LT tangents at $U \equiv \mathbb{1}$
  & expected rank $5$ within $10^{-10}$ \\
8 & Gauge-cleaning projector
  & $P_\perp = \mathbb{1} - V_{\rm gauge}V_{\rm gauge}^{\!\top}$
  & $P_\perp^2 = P_\perp$ within $10^{-14}$ \\
9 & Gram positivity
  & $5\times 5$ $G_{ab} = \langle t_a, t_b\rangle$
  & $\det G > 0$ and $\lambda_{\min}(G) > 0$ \\
\bottomrule
\end{tabular}
\caption{Infrastructure verification suite. All nine tests return pass
on every production row used in the main text.}
\label{tab:app-infra}
\end{table}

The current repository state used for the production runs reports this
suite as $9/9$; older development logs with a different count reflect
superseded draft-stage bookkeeping and are not part of the manuscript
baseline.

Tests~1--3 verify the lattice geometry and gauge-field assembly.
Tests~4--6 verify consistency of the three operator constructors with
the sign and regulator conventions of Appendix~\ref{app:conventions}.
Tests~7--9 verify the Jacobian infrastructure: the left-trivialised
tangents recover the expected five-dimensional tangent space at a
trivial gauge background, the gauge-cleaning projector is idempotent
within numerical precision, and the $5\times 5$ Gram assembled from the
cleaned tangents is strictly positive-definite. On every production row
used in the main text, the suite returns $9/9$.

The full implementation, including the numerical tolerances entering
each criterion, is available in the public repository referenced in the
data-availability statement of the main text.



\section{$\Icl$-degraded rows}
\label{app:icl-degraded}

The primary production rows A, B, C of Section~\ref{sec:finite-xi} span
$\Icl \in [0.983,\,1.032]$, so that the clover discretization tracks the
continuum target $\Icl = 1$ at the few-percent level. Additional
exploratory rows at larger $\rho/L$ or smaller $V$ yield systematically
larger $\Icl$-values. They are retained here only as a systematics
reference for the dependence of $\delta_\xi \log\zeta$ on $\Icl$; they
are \emph{not} used in the assembled result of
Section~\ref{sec:assembled}.

Table~\ref{tab:app-icl-degraded} lists the four reference rows and their
Simpson values, together with the ratio $\delta_\xi \log\zeta / \epsilon$
that renders the trend below dimensionless.

\begin{table}[h]
\centering
\begin{tabular}{lccccc}
\toprule
Row & Dimensions & $\rho$ & $\Icl$ & $\delta_\xi \log\zeta$ &
$\delta_\xi \log\zeta / \epsilon$ \\
\midrule
D1 & $8^3 \times 8$    & $3.0$ & $1.201$ & $-2.240 \pm 0.168$ & $-24.9$ \\
D2 & $12^3 \times 8$   & $3.0$ & $1.068$ & $-2.095 \pm 0.174$ & $-23.3$ \\
D3 & $18^3 \times 12$  & $4.5$ & $1.218$ & $-2.802 \pm 0.168$ & $-31.1$ \\
D4 & $24^3 \times 16$  & $6.0$ & $1.329$ & $-3.600 \pm 0.384$ & $-40.0$ \\
\bottomrule
\end{tabular}
\caption{Rows with $\Icl$ outside the primary production range. All
quoted $\delta_\xi \log\zeta$ values are Simpson measurements in the
same convention used in the main text. These rows are retained only as
a systematics reference and are not used in the assembled result.}
\label{tab:app-icl-degraded}
\end{table}

The reference rows should be read only as a one-parameter trend study
in $\Icl$, not as an alternative measurement set: they probe how the
branch response moves once the caloron is visibly under-resolved.

A weighted linear regression through the four reference rows, using
the statistical uncertainties quoted in
Table~\ref{tab:app-icl-degraded}, gives
\begin{equation}
\frac{\delta_\xi \log\zeta}{\epsilon}
\;\approx\;
-\,18.7 \;-\; 49.5\,(\Icl - 1),
\label{eq:app-icl-correlation}
\end{equation}
over the range $\Icl \in [1.07,\,1.33]$ spanned by
Table~\ref{tab:app-icl-degraded}. The fit has
$\chi^2/\mathrm{dof} = 3.3$, dominated by Row~D1 at a
$+2.0\sigma$ residual; the other three rows sit within $1.2\sigma$ of
the fit line. To account conservatively for this scatter, the
extrapolation uncertainty below is scale-inflated by
$\sqrt{\chi^2/\mathrm{dof}} \simeq 1.82$.

\begin{figure}[h]
\centering
\includegraphics[width=0.60\textwidth]{fig_icl_correlation.pdf}
\caption{The four $\Icl$-degraded reference rows (D1--D4) plotted as
$\delta_\xi \log\zeta / \epsilon$ versus $\Icl - 1$, with the
weighted linear-regression guide~\eqref{eq:app-icl-correlation}
shown as a solid line ($\chi^2/\mathrm{dof} = 3.3$, dominated by
Row~D1 at $+2.0\sigma$ residual). The dashed vertical line marks
$\Icl = 1$; the red square at that point shows the extrapolated
value with $\sqrt{\chi^2/\mathrm{dof}}$-inflated uncertainty,
statistically consistent with the primary three-row weighted result
$-1.669 \pm 0.125$ of Eq.~\eqref{eq:bracket-result}.}
\label{fig:app-icl-correlation}
\end{figure}

Extrapolated to $\Icl = 1$, \eqref{eq:app-icl-correlation} gives
\begin{equation}
\delta_\xi \log\zeta\bigr|_{\Icl = 1}
\;=\; -1.68 \pm 0.44,
\label{eq:app-icl-extrap}
\end{equation}
which agrees with the primary three-row weighted
result~\eqref{eq:bracket-result}, $-1.669 \pm 0.125$, at the
$0.03\sigma$ level. This independent extrapolation from rows
visibly outside the admissibility window provides an additional
cross-check on the production value of the bracket: the
$\Icl$-dependence inferred from rows three to ten percent away from
the continuum target is consistent with the in-window result within
statistical error.


The physical origin of $\Icl > 1$ in the reference rows is the standard
under-resolution of the caloron core at large $\rho/L$: the clover
field-strength tensor systematically overestimates
$\Tr(F_{\mu\nu}F^{\mu\nu})$ when the caloron profile varies appreciably
across a single plaquette. The effect is a lattice artefact of the
discretization, not a property of the continuum background, and is
removed by working at the smaller $\rho/L$ values of the production
rows. Rows~D1 and~D2 illustrate the opposite limit of the same effect
at small $V$ and small $\rho$: the $\Icl$-deviation there is driven by
finite-volume rather than finite-spacing artefacts, and both combine
into the single-parameter trend~\eqref{eq:app-icl-correlation}.

The production campaign of Section~\ref{sec:finite-xi} uses only rows
with $\Icl$ close to unity, for which the left-hand side of
\eqref{eq:app-icl-correlation} is insensitive to the remaining
$\Icl$-dependence and the lattice artefact is subdominant to the
statistical error.



\section{Bracket error budget}
\label{app:error-budget}

The three-lattice weighted bracket reported in Section~\ref{sec:finite-xi},
\begin{equation}
\delta_\xi \log\zeta \;=\; -1.669 \pm 0.125,
\qquad
\Bobs \;=\; 0.811 \pm 0.023,
\label{eq:app-bracket-weighted}
\end{equation}
has its stated uncertainty dominated by the Hutchinson trace estimator on
the individual production rows. This appendix decomposes the $\pm 0.023$
uncertainty on $\Bobs$ into its contributing components so that a reader
can reproduce the weighted error directly from the row-level entries of
Table~\ref{tab:production-results}.

\subsection{Per-row statistical uncertainty}
\label{app:err-row-statistical}

On each production row, the quoted quantity $\delta_\xi \log\zeta$ is the
fully assembled three-node Simpson estimator built from the paired-noise
node integrands $h(\xi_j)$:
\begin{equation}
\delta_\xi \log\zeta \;=\;
\frac{\epsilon}{6}\Bigl[h(1) + 4\,h(1+\epsilon/2) + h(1+\epsilon)\Bigr]
\;+\; J_{\rm anchor},
\label{eq:app-simpson}
\end{equation}
where $J_{\rm anchor} = \Jphys$ is the anchor-Jacobian contribution
of~\eqref{eq:Janchor-equals-Jphys}. Because $J_{\rm anchor}$ is
$\xi$-independent, it contributes to the mean of $\delta_\xi \log\zeta$
but not to its Hutchinson variance; the variance analysis below
therefore treats only the $\xi$-dependent Simpson combination.
Under the paired-Hutchinson construction of
Appendix~\ref{app:conventions}, the three node integrands
$h(\xi_0), h(\xi_1), h(\xi_2)$ are computed against a common
Rademacher realisation, so that their cross-covariances are
non-trivial and largely positive. The per-row statistical
uncertainty on $\delta_\xi \log\zeta$ is therefore the standard
error of the full Simpson sample,
\begin{equation}
\sigma_{\rm stat}(\delta_\xi \log\zeta)
\;=\;
\frac{1}{\sqrt{N_S}}\,
\mathrm{sd}\!\left[
\frac{\epsilon}{6}\bigl(
H_{0,s} + 4\,H_{1,s} + H_{2,s}
\bigr)
\right]_{s = 1,\dots,N_S},
\label{eq:app-sigma-stat}
\end{equation}
where $H_{j,s}$ denotes the $s$-th Hutchinson realisation of the
integrand at node $\xi_j$ and $\mathrm{sd}[\cdot]$ is the
sample standard deviation over paired samples. Computing
\eqref{eq:app-sigma-stat} directly on the production samples
reproduces the row-level errors $0.186$, $0.213$, and $0.275$ of
Table~\ref{tab:production-results} within rounding. The gain of
the paired construction is not a formal one-term factor in
\eqref{eq:app-sigma-stat} but the fact that the positive
cross-covariances between $H_{0,s}$, $H_{1,s}$, and $H_{2,s}$
reduce $\mathrm{sd}[\cdot]$ by roughly an order of magnitude
relative to the uncorrelated case.

\subsection{Weighted combination}
\label{app:err-weighted}

The three production rows are combined by inverse-variance weighting,
\begin{equation}
\delta_\xi \log\zeta
\;=\;
\frac{\sum_i w_i\, (\delta_\xi \log\zeta)_i}{\sum_i w_i},
\qquad
w_i = \sigma_i^{-2},
\label{eq:app-weighted-mean}
\end{equation}
with weights $(w_{\rm A},\,w_{\rm B},\,w_{\rm C}) = (28.9,\,22.0,\,13.2)$
and weighted standard error
\begin{equation}
\sigma_{\rm w}(\delta_\xi \log\zeta)
\;=\;
\Bigl(\sum_i w_i\Bigr)^{\!-1/2}
\;=\; 0.125.
\label{eq:app-sigma-w}
\end{equation}
These weighted row-level uncertainties are the only inputs needed for
the bracket error budget. Other factors entering the assembled
cosmological constant, such as $\Jphys$, belong to the global error
budget of Section~\ref{sec:assembled} and not to the bracket budget
itself.

The three row residuals relative to the weighted mean are $0.009$,
$-0.026$, and $0.023$, well within the row-level errors and consistent
with the pairwise tensions of Section~\ref{sec:continuum}.

\subsection{Bracket nonlinearity}
\label{app:err-bracket-nonlinearity}

In the sign convention used throughout the paper, the weighted branch
response is negative,
\[
\delta_\xi \log\zeta \;=\; -1.669 \pm 0.125,
\]
and the bracket is therefore
\begin{equation}
\Bobs \;=\; 1 - \exp\!\bigl(\delta_\xi \log\zeta\bigr).
\label{eq:app-bracket-map}
\end{equation}
Equivalently, if one introduces the positive magnitude
$\Delta_\zeta \equiv -\delta_\xi \log\zeta > 0$, then
$\Bobs = 1 - e^{-\Delta_\zeta}$. We keep the former convention here
because it matches the main-text variable.

The bracket uncertainty follows by first-order propagation:
\begin{equation}
\sigma(\Bobs)
\;=\;
\left|\frac{\partial \Bobs}{\partial(\delta_\xi\log\zeta)}\right|
\,\sigma_{\rm w}(\delta_\xi \log\zeta)
\;=\;
e^{\delta_\xi \log\zeta}\,\sigma_{\rm w}(\delta_\xi \log\zeta)
\;=\; 0.188 \times 0.125
\;=\; 0.023.
\label{eq:app-sigma-bracket}
\end{equation}
This reproduces the quoted weighted uncertainty on $\Bobs$. At the
current precision the nonlinearity of the map
$\Bobs = 1 - e^{\delta_\xi\log\zeta}$ is negligible for the uncertainty
propagation. The leading second-order correction to the mean is
\begin{equation}
\frac{1}{2}\,e^{\delta_\xi \log\zeta}\,
\sigma_{\rm w}^2(\delta_\xi \log\zeta)
\;\simeq\; 1.5\times 10^{-3},
\label{eq:app-bracket-second-order}
\end{equation}
which affects only the third decimal place of $\Bobs$.

\subsection{Sub-dominant systematics}
\label{app:err-sub-dominant}

The remaining bracket-level systematics contribute below the
statistical uncertainty on $\Bobs$ and are not propagated
into~\eqref{eq:app-bracket-weighted}.
Table~\ref{tab:app-error-budget} collects them for completeness. The
FD-step entry refers to the convergence of the Simpson node values
$h(\xi_j)$ under halving of the finite-difference step used in the
tangent construction. The $\Icl$-trend entry refers to the variation
induced across the production range $\Icl \in [0.983, 1.032]$, which is
already captured in the row-to-row spread feeding the inverse-variance
combination \eqref{eq:app-weighted-mean}.

\begin{table}[h]
\centering
\begin{tabular}{llc}
\toprule
Source & Treatment & Contribution to $\sigma(\Bobs)$ \\
\midrule
Hutchinson statistics, three-row weighted
  & included in~\eqref{eq:app-sigma-bracket}
  & $2.3 \times 10^{-2}$ \\
Paired-difference covariance
  & included in the per-row $\sigma_i$
  & (subsumed) \\
FD step in tangent construction
  & convergence $\delta \to \delta/2$
  & $< 10^{-4}$ \\
CG residual ($10^{-10}$)
  & solver tolerance in projected CG
  & $< 10^{-6}$ \\
$\Icl$-trend across the production range
  & already included in the row-to-row spread
  & (subsumed) \\
Simpson truncation ($\epsilon = 0.09$)
  & beyond three-point quadrature
  & $< 10^{-3}$ \\
\bottomrule
\end{tabular}
\caption{Contributions to the bracket uncertainty and their treatment.
Entries listed as ``subsumed'' are already included in either the
per-row $\sigma_i$ or the weighted combination.}
\label{tab:app-error-budget}
\end{table}

The bracket uncertainty reported in~\eqref{eq:app-bracket-weighted} is
therefore $\sigma(\Bobs) = 0.023$, dominated by the Hutchinson
statistics on the three production rows, with all other systematic
contributions subdominant at the current precision.


\section{$S_{\rm inst}$-sensitivity of the assembled result}
\label{app:S-sensitivity}

The master formula~\eqref{eq:master} separates the assembled
cosmological constant into an exponential that depends only on
$S_{\rm inst}$ and a lattice-determined prefactor that does not:
\begin{equation}
\frac{\Lobs}{\MP^4} \;=\; \mathcal{K}\cdot e^{-S_{\rm inst}},
\qquad
\mathcal{K}
\;\equiv\;
\Jphys\,\bigl[\zeta(0) - \zeta(X_*)\bigr]\,\cos\Theta_\gamma
\;=\;
\Jphys\,\zeta(0)\,\Bobs\,\cos\Theta_\gamma,
\label{eq:app-S-factorisation}
\end{equation}
with $S_{\rm inst} = 8\pi^2 a_C \Ilat$. This appendix records how
$\Lobs/\MP^4$ responds to variations of $S_{\rm inst}$ at fixed
$\mathcal{K}$.

The factors entering $\mathcal{K}$ are fixed by the lattice computation
of the Jacobian, the phase, and the branch-response bracket, together
with the normalization convention used in
Section~\ref{sec:assembled}. For the present appendix it is convenient
to work with the central assembled prefactor multiplying $10^{-120}$,
namely $0.81$, and to define $S^\star$ by
\begin{equation}
0.81 \times 10^{-120} \;=\; \mathcal{K}\,e^{-S^\star}.
\label{eq:app-Sstar-def}
\end{equation}
Numerically $S^\star \simeq 276.3$. The decimal place here reflects
only the normalization convention of the table; the main text
continues to state the hidden-sector input as
$S_{\rm inst} \simeq 276$. This value is only the normalization point
used for the sensitivity table below; it is not a separate
hidden-sector input.

Differentiating~\eqref{eq:app-S-factorisation},
\begin{equation}
\frac{\partial \ln\bigl(\Lobs/\MP^4\bigr)}{\partial S_{\rm inst}}
\;=\; -\,1,
\label{eq:app-S-sensitivity}
\end{equation}
so that a unit shift in $S_{\rm inst}$ moves $\Lobs/\MP^4$ by one
$e$-fold. Equivalently,
\[
\frac{\Lobs/\MP^4}{0.81\times 10^{-120}}
\;=\;
e^{-(S_{\rm inst} - S^\star)},
\]
so the table below depends only on the shift
$S_{\rm inst} - S^\star$. Table~\ref{tab:app-S-sensitivity} evaluates
the response of the assembled result across a range of approximately thirty $e$-folds around $S^\star$.

\begin{table}[h]
\centering
\begin{tabular}{cccl}
\toprule
$S_{\rm inst}$ &
$e^{-(S_{\rm inst} - S^\star)}$ &
$\Lobs/\MP^4 \,/\, 10^{-120}$ &
Deviation from observed \\
\midrule
$260.3$ & $8.9 \times 10^{+6}$ & $7.2 \times 10^{+6}$ & $+7$ orders \\
$270.3$ & $4.0 \times 10^{+2}$ & $3.3 \times 10^{+2}$ & $+2.6$ orders \\
$276.3$ & $1.0$                & $0.81$               & central \\
$280.3$ & $1.8 \times 10^{-2}$ & $1.5 \times 10^{-2}$ & $-1.7$ orders \\
$290.3$ & $8.3 \times 10^{-7}$ & $6.7 \times 10^{-7}$ & $-6$ orders \\
\bottomrule
\end{tabular}
\caption{Response of the assembled cosmological constant to variations
of $S_{\rm inst}$. The central row is defined by
$\mathcal{K}\,e^{-S^\star} = 0.81 \times 10^{-120}$ at
$\mathcal{K} = 0.81$, giving $S^\star \simeq 276.3$. Shifts of
$S_{\rm inst}$ by a few units change $\Lobs/\MP^4$ by orders of
magnitude.}
\label{tab:app-S-sensitivity}
\end{table}

Two observations follow.

First, the $\mathcal{O}(1)$ agreement between the assembled result and
the observed cosmological constant is a quantitative match, not a
tautological one. A shift of $S_{\rm inst}$ by $\pm 5$ would change
$\Lobs/\MP^4$ by roughly two orders of magnitude, far outside the $3\%$
precision of the lattice prefactor. The narrow window within which the
match holds at the prefactor level corresponds, by
\eqref{eq:app-S-sensitivity}, to the fractional prefactor uncertainty
propagated through the logarithmic derivative.

Second, if one holds the assembled prefactor fixed and asks how
sharply the match depends on $S_{\rm inst}$, the allowed window is
narrow:
\begin{equation}
|\Delta S_{\rm inst}|
\;\lesssim\;
\frac{\sigma(\mathcal{K})}{\mathcal{K}}
\;\simeq\; 0.04 .
\label{eq:app-S-precision}
\end{equation}
This should be read only as a fixed-framework sensitivity statement.
It is not an independent measurement of hidden-sector dynamics, and it
does not reduce the conditional status of the main result.


\bibliographystyle{iopart-num}
\bibliography{bib}

\end{document}